\part{Patrón MVC y Servicios Web}

\chapter{El patrón Modelo-Vista-Controlador en aplicaciones web modernas}
\label{cap:semana14}

\epigraph{«MVC no es un framework: es una forma de pensar el software
que separa lo que el usuario ve, lo que el sistema sabe, y la lógica
que conecta ambos mundos.»}{--- \textit{Adaptado de Trygve Reenskaug}}

\section*{Resultado de aprendizaje de la semana}
Al concluir la semana 14, el estudiante diseña aplicaciones web bajo
el patrón Modelo-Vista-Controlador, comprende la trayectoria histórica
desde el MVC original de Smalltalk-80 hasta sus variantes modernas
(server-side MVC vs API REST + SPA), y construye una aplicación
funcional con un \emph{framework} MVC contemporáneo, depurándola paso
a paso en IntelliJ IDEA Ultimate.

\section{Historia del patrón MVC}

MVC es uno de los patrones de diseño más citados en la
historia del software. Comprender su historia ---desde Smalltalk-76
hasta los frameworks modernos--- ayuda a entender por qué hay tantas
variantes y por qué sigue siendo central en 2026.


\begin{cajahistoria}{Reenskaug --- Smalltalk-80 y el nacimiento de un patrón fundamental}
El patrón MVC fue concebido por el ingeniero noruego \textbf{Trygve
Reenskaug} durante una estancia sabática en el Xerox PARC en
diciembre de 1978 mientras trabajaba en la primera versión de
Smalltalk. Reenskaug describió MVC como una solución al problema
fundamental de los sistemas interactivos: cómo permitir que múltiples
vistas reflejen el mismo modelo de datos sin acoplarlas entre sí.

En su nota original \emph{«Models-Views-Controllers»} (10 de
diciembre de 1979), Reenskaug definía:

\begin{itemize}
\item \textbf{Modelo}: la representación del conocimiento sobre el
problema; lo que la aplicación realmente \emph{sabe}.
\item \textbf{Vista}: una representación visual del modelo, no el
modelo mismo.
\item \textbf{Controlador}: el enlace entre el usuario y el sistema,
que interpreta las acciones del usuario y las traduce en cambios al
modelo.
\end{itemize}

MVC fue implementado plenamente en Smalltalk-80 (1980) y durante
quince años permaneció en relativa oscuridad académica. Su
revitalización vino con dos hitos:

\begin{itemize}
\item \textbf{NeXTSTEP / Cocoa} (1989+): Apple adoptó MVC como
patrón estructural de la programación nativa de Mac y, posteriormente,
de iOS.
\item \textbf{Ruby on Rails} (David Heinemeier Hansson, 2004):
implementó una variante \emph{server-side} de MVC para aplicaciones
web. Su éxito masivo inspiró Django (Python, 2005), Symfony
(PHP, 2005), CodeIgniter (PHP, 2006), CakePHP (PHP, 2005), Laravel
(PHP, 2011), ASP.NET MVC (Microsoft, 2009) y Spring MVC (2004).
\end{itemize}

A partir de 2014 emergieron las arquitecturas \textbf{API REST + SPA}
que desplazaron parcialmente el MVC server-side: el servidor expone
JSON, y un cliente Angular/React/Vue se encarga de la vista. Sin
embargo, el patrón MVC sigue siendo fundamental: incluso las SPAs
implementan internamente versiones del patrón (MVVM, Flux, Redux,
Signal-based architectures) que son herederas conceptuales del MVC
original.
\end{cajahistoria}

\subsection{Tabla evolutiva del MVC en frameworks web}

Distintos frameworks implementan MVC con variantes propias.
La tabla siguiente recorre los más influyentes en orden cronológico.


\begin{table}[htbp]
\centering
\caption{Hitos en la historia del MVC en frameworks web.}
\footnotesize
\begin{tabularx}{\textwidth}{lllX}
\toprule
\textbf{Año} & \textbf{Framework} & \textbf{Lenguaje} & \textbf{Aporte}\\
\midrule
1979 & ---           & Smalltalk & Reenskaug formaliza el patrón.\\
1980 & Smalltalk-80  & Smalltalk & Primera implementación completa.\\
1989 & NeXTSTEP      & Objective-C & MVC en aplicaciones nativas de escritorio.\\
1996 & Apache Struts (precursor) & Java & MVC web temprano (Java Servlets).\\
2002 & Spring        & Java & DI + MVC posteriormente (Spring MVC, 2004).\\
2004 & Ruby on Rails & Ruby & Convención sobre configuración; MVC mainstream.\\
2005 & Django        & Python & Modelo MTV (Model-Template-View, variante MVC).\\
2005 & Symfony       & PHP   & MVC con bundles y dependencias.\\
2009 & ASP.NET MVC   & C\#   & Microsoft adopta MVC web.\\
2011 & Laravel       & PHP   & MVC con Eloquent ORM y Blade templates.\\
2014 & Spring Boot   & Java  & Convención sobre Spring MVC, simplificación radical.\\
2014 & Angular 2     & TypeScript & Frontend MVC (versión MVVM).\\
2024 & Laravel 11    & PHP   & Estructura simplificada, Volt, Folio.\\
2024 & Spring Boot 4 & Java & Soporte HTTP virtual threads en Tomcat.\\
2024 & ASP.NET Core 10 & C\#  & Razor Pages + MVC + Web API unificados.\\
\bottomrule
\end{tabularx}
\end{table}

\section{Anatomía del patrón MVC en la web}

MVC en aplicaciones web tiene una anatomía específica: el
modelo no es solo una clase, sino una capa entera; la vista no es
una pantalla, sino una plantilla; el controlador no es un widget,
sino un endpoint HTTP.


\subsection{Adaptación del MVC original al contexto HTTP}

El MVC de Reenskaug fue concebido para aplicaciones de escritorio
con observadores reactivos: la vista se suscribía al modelo, que la
notificaba en cada cambio. En la web esto no es posible directamente
(HTTP es \emph{request/response}), por lo que el MVC web introduce
adaptaciones:

\begin{itemize}
\item No hay observadores reactivos: la vista se renderiza en cada
petición.
\item Hay un \textbf{front controller} (DispatcherServlet en Spring,
\texttt{index.php} en Laravel, \texttt{Routes} en ASP.NET Core) que
recibe TODAS las peticiones y las enruta al controller adecuado.
\item Las URLs se mapean a métodos de controller con sistemas de
\emph{routing} declarativos.
\item Las vistas se generan típicamente en el servidor con
\emph{template engines} (Thymeleaf, Blade, Razor, Twig).
\end{itemize}

\subsection{Las tres capas en detalle}

Cada una de las tres capas de MVC tiene responsabilidades
concretas. La descripción siguiente las detalla.


\subsubsection{El Modelo}

El Modelo encapsula:
\begin{itemize}
\item \textbf{Entidades de dominio}: objetos persistentes con
identidad propia (Producto, Pedido, Usuario).
\item \textbf{Servicios de aplicación}: orquestan casos de uso
(crear pedido, calcular descuento, enviar notificación).
\item \textbf{Repositorios}: encapsulan el acceso a la persistencia.
\item \textbf{Reglas de negocio}: invariantes del dominio
(un usuario menor de edad no puede comprar bebidas alcohólicas).
\item \textbf{DTOs} (Data Transfer Objects): contenedores planos
para transferir datos entre capas o servicios.
\end{itemize}

\textbf{Anti-patrón clásico:} el «Anemic Domain Model» (Fowler) en
el que las entidades son meros contenedores sin lógica y toda la
inteligencia está en servicios. En MVC profesional las entidades
tienen métodos significativos.

\subsubsection{La Vista}

La Vista es la capa de presentación. En aplicaciones web tradicionales
suele ser una plantilla HTML con marcadores de sustitución:

\begin{itemize}
\item \textbf{Thymeleaf} (Java/Spring): plantillas HTML que son HTML
válido incluso sin procesar (preview-friendly).
\item \textbf{Blade} (PHP/Laravel): \texttt{@if}, \texttt{@foreach},
\texttt{\{\{ \$variable \}\}}.
\item \textbf{Razor} (C\#/ASP.NET): mezcla \texttt{@} para C\# inline.
\item \textbf{Twig} (PHP/Symfony): sintaxis estilo Django,
\texttt{\{\% if \%\} \{\{ var \}\}}.
\item \textbf{Mustache, Handlebars} (cualquier lenguaje): \emph{logic-less}.
\end{itemize}

En arquitecturas API REST + SPA, la Vista se desplaza al cliente
(Angular, React, Vue), y el servidor solo expone JSON.

\subsubsection{El Controlador}

El Controller es la capa que recibe peticiones HTTP, invoca al Modelo
y selecciona la Vista (o el formato de respuesta JSON). Sus
responsabilidades:

\begin{itemize}
\item \emph{Routing}: mapear URLs a métodos.
\item \emph{Binding}: convertir parámetros HTTP en objetos del modelo.
\item \emph{Validation}: validar entrada (a menudo delegada).
\item \emph{Orquestación}: llamar a servicios del modelo.
\item \emph{Selección de respuesta}: vista a renderizar o DTO a
serializar.
\item \emph{Manejo de errores}: convertir excepciones en respuestas
HTTP apropiadas.
\end{itemize}

\textbf{Regla de oro:} el controller debe ser \emph{delgado} (\emph{thin
controller}). Toda lógica de negocio significativa va en el modelo
(servicios). Un controller con 200 líneas es señal de mal diseño.

\section{Ciclo de vida de una petición MVC en Spring Boot}

Una petición HTTP atraviesa una secuencia precisa de
componentes dentro de Spring Boot. Conocer ese flujo permite
depurar problemas y optimizar puntos críticos.


\begin{figure}[htbp]
\centering
\resizebox{\textwidth}{!}{%
\begin{tikzpicture}[
  every node/.style={font=\scriptsize,align=center},
  caja/.style={rectangle,draw=uteqverde,thick,minimum width=2.4cm,
               minimum height=0.8cm,fill=uteqverde!10,rounded corners=2pt},
  flecha/.style={->,>=stealth,thick}]

  \node[caja] (cli) {1. Cliente\\(Navegador)};
  \node[caja,below=0.5cm of cli] (ds) {2. Dispatcher\\Servlet};
  \node[caja,below=0.5cm of ds] (hm) {3. Handler\\Mapping};
  \node[caja,below=0.5cm of hm] (ctr) {4. Controller\\@RequestMapping};
  \node[caja,right=1.5cm of ctr] (svc) {5. Service\\(reglas negocio)};
  \node[caja,right=1.5cm of svc] (rp) {6. Repository\\(JPA)};
  \node[cylinder,draw=uteqverde,fill=uteqverde!15,thick,
        shape border rotate=90,aspect=0.3,minimum width=1.2cm,
        minimum height=1.2cm,below=0.5cm of rp] (db) {BD};
  \node[caja,left=1.5cm of ctr] (vr) {8. View\\Resolver};
  \node[caja,left=1.5cm of vr] (vw) {9. Vista\\(Thymeleaf)};

  \draw[flecha] (cli) -- node[right]{HTTP Request} (ds);
  \draw[flecha] (ds) -- (hm);
  \draw[flecha] (hm) -- (ctr);
  \draw[flecha] (ctr) -- (svc);
  \draw[flecha] (svc) -- (rp);
  \draw[flecha] (rp) -- (db);
  \draw[flecha,dashed] (db) -- node[right,font=\tiny]{rows} (rp);
  \draw[flecha,dashed] (rp) -- (svc);
  \draw[flecha,dashed] (svc) -- (ctr);
  \draw[flecha] (ctr) -- node[above,font=\tiny]{ModelAndView} (vr);
  \draw[flecha] (vr) -- (vw);
  \draw[flecha,dashed] (vw.north) to[bend left=60]
        node[above,font=\tiny]{HTML} (cli.west);

\end{tikzpicture}}
\caption{Ciclo de vida completo de una petición HTTP en Spring MVC. La
ida (línea sólida): Dispatcher $\to$ HandlerMapping $\to$ Controller
$\to$ Service $\to$ Repository $\to$ BD. El regreso (línea punteada):
los datos suben hasta el Controller, que devuelve un
\texttt{ModelAndView}; el ViewResolver elige la plantilla y el HTML
final llega al cliente.}
\end{figure}

\section{Server-side MVC vs SPA + API REST}

En 2026 hay dos arquitecturas dominantes para aplicaciones
web: MVC \emph{server-side} clásico (Thymeleaf/Razor/Blade) y SPA +
API REST (Angular/React + Spring/ASP.NET). La tabla siguiente las
contrasta con sus trade-offs.


\begin{cajacompara}{MVC server-side tradicional vs API REST + SPA Angular/React}
\centering\footnotesize
\begin{tabularx}{\textwidth}{lXX}
\toprule
\textbf{Criterio} & \textbf{MVC server-side} & \textbf{API REST + SPA}\\
\midrule
Ubicación de la vista & Servidor (Thymeleaf, Blade, Razor) & Cliente (Angular, React, Vue)\\
Renderizado inicial & Rápido (HTML completo del primer hit) & Lento (descarga JS, luego renderiza)\\
Tráfico tras la primera carga & Cada navegación = HTML completo & JSON pequeño\\
SEO & Excelente nativo & Requiere SSR/SSG\\
Complejidad del frontend & Baja (HTML + algo de JS) & Alta (framework completo)\\
Reactividad & Limitada (post-back o JS extra) & Nativa\\
Equipo & Un equipo full-stack & Frontend y backend separados\\
Mejor para & CMS, intranet, sitios contenido & SaaS, dashboards, apps tipo Gmail\\
\bottomrule
\end{tabularx}
\end{cajacompara}

\subsection{¿Cuál usar en el PFC?}

El PFC de esta asignatura adopta la combinación \textbf{API REST
(Spring Boot) + SPA Angular}. Es la apuesta arquitectónica moderna y
es la que el sílabo prescribe. Sin embargo, comprender ambos
enfoques es esencial: muchos sistemas reales coexisten (un panel
admin con Razor + una API para móvil + un portal público con SSR).

\section{Implementación canónica: MVC server-side completo}

El ejemplo siguiente implementa un CRUD MVC server-side
completo con Spring Boot 4 y Thymeleaf: el estudiante lo construye
paso a paso y obtiene una aplicación funcional ejecutable.


\begin{cajaejemplo}{Ejemplo 14.1 --- CRUD MVC clásico en Spring Boot 4 + Thymeleaf}


\textbf{Enunciado.} Construir una aplicación MVC server-side de gestión de tareas (to-do) con Spring Boot 4, Thymeleaf y H2. Las cuatro operaciones CRUD funcionarán end-to-end.

\textbf{Estructura del proyecto (Listado~\ref{lst:14.1}):}
\begin{lstlisting}[style=shell,numbers=none,caption={Ejemplo de Shell},label=lst:14.1]
src/main/java/ec/edu/uteq/mvc/
  controller/ProductoController.java   <- Controller
  service/ProductoService.java         <- Servicio (modelo)
  repository/ProductoRepository.java   <- Repositorio (modelo)
  model/Producto.java                  <- Entidad de dominio
  dto/ProductoForm.java                <- DTO con validacion
src/main/resources/templates/productos/ <- Vistas
  lista.html
  formulario.html
\end{lstlisting}

\textbf{Entidad \texttt{Producto}:}
\begin{lstlisting}[style=java,caption={Producto.java},label=lst:14.2]
package ec.edu.uteq.mvc.model; // paquete del archivo

import jakarta.persistence.*; // importacion de clase externa
import lombok.*; // importacion de clase externa
import java.math.BigDecimal; // importacion de clase externa

@Entity @Table(name = "productos") // entidad JPA (mapea a tabla)
@Getter @Setter @NoArgsConstructor @AllArgsConstructor
public class Producto { // clase publica
    @Id @GeneratedValue(strategy = GenerationType.IDENTITY) // campo clave primaria
    private Long id; // campo privado
    private String nombre; // campo privado
    private BigDecimal precio; // campo privado
    private Integer stock; // campo privado

    /** Logica de dominio: el modelo NO es anemico. */
    public boolean estaDisponible() { // metodo publico
        return stock != null && stock > 0; // retorna
    }

    public BigDecimal valorInventario() { // metodo publico
        // condicional
        // condicional
        if (precio == null || stock == null) return BigDecimal.ZERO;
        return precio.multiply(BigDecimal.valueOf(stock)); // retorna
    }
}
\end{lstlisting}

\textbf{DTO con validación (Listado~\ref{lst:14.3}):}
\begin{lstlisting}[style=java,caption={ProductoForm.java},label=lst:14.3]
package ec.edu.uteq.mvc.dto; // paquete del archivo

import jakarta.validation.constraints.*; // importacion de clase externa
import lombok.*; // importacion de clase externa
import java.math.BigDecimal; // importacion de clase externa

@Getter @Setter @NoArgsConstructor @AllArgsConstructor
public class ProductoForm { // clase publica
    private Long id; // campo privado

    @NotBlank(message = "El nombre es obligatorio") // validacion: no nulo y no vacio
    @Size(min = 3, max = 120) // validacion: longitud min/max
    private String nombre; // campo privado

    @NotNull @DecimalMin("0.00") @DecimalMax("99999.99") // validacion: no nulo
    private BigDecimal precio; // campo privado

    @NotNull @Min(0) // validacion: no nulo
    private Integer stock; // campo privado
}
\end{lstlisting}

\textbf{Repository (Spring Data) (Listado~\ref{lst:14.4}):}
\begin{lstlisting}[style=java,caption={ProductoRepository.java},label=lst:14.4]
package ec.edu.uteq.mvc.repository; // paquete del archivo

import ec.edu.uteq.mvc.model.Producto; // importacion de clase externa
// importacion de clase externa
// importacion de clase externa
import org.springframework.data.jpa.repository.JpaRepository;

public interface ProductoRepository // declaracion de interfaz
    extends JpaRepository<Producto, Long> {}
\end{lstlisting}

\textbf{Service (Modelo) (Listado~\ref{lst:14.5}):}
\begin{lstlisting}[style=java,caption={ProductoService.java},label=lst:14.5]
package ec.edu.uteq.mvc.service; // paquete del archivo

import ec.edu.uteq.mvc.model.Producto; // importacion de clase externa
import ec.edu.uteq.mvc.repository.ProductoRepository; // importacion de clase externa
import lombok.RequiredArgsConstructor; // importacion de clase externa
import org.springframework.stereotype.Service; // importacion de clase externa
// importacion de clase externa
// importacion de clase externa
import org.springframework.transaction.annotation.Transactional;
import java.util.*; // importacion de clase externa

@Service // capa de servicio (logica de negocio)
@RequiredArgsConstructor
public class ProductoService { // clase publica

    private final ProductoRepository repo; // inmutable

    @Transactional(readOnly = true) // envuelve el metodo en transaccion
    // metodo publico
    // metodo publico
    public List<Producto> listar() { return repo.findAll(); }

    @Transactional(readOnly = true) // envuelve el metodo en transaccion
    public Producto buscar(Long id) { // metodo publico
        return repo.findById(id).orElseThrow( // retorna
            () -> new NoSuchElementException("Producto " + id));
    }

    @Transactional // envuelve el metodo en transaccion
    // metodo publico
    // metodo publico
    public Producto guardar(Producto p) { return repo.save(p); }

    @Transactional // envuelve el metodo en transaccion
    public void eliminar(Long id) { repo.deleteById(id); } // metodo publico
}
\end{lstlisting}

\textbf{Controller (Listado~\ref{lst:14.6}):}
\begin{lstlisting}[style=java,caption={ProductoController.java},label=lst:14.6]
package ec.edu.uteq.mvc.controller; // paquete del archivo

import ec.edu.uteq.mvc.dto.ProductoForm; // importacion de clase externa
import ec.edu.uteq.mvc.model.Producto; // importacion de clase externa
import ec.edu.uteq.mvc.service.ProductoService; // importacion de clase externa
import jakarta.validation.Valid; // importacion de clase externa
import lombok.RequiredArgsConstructor; // importacion de clase externa
import org.springframework.beans.BeanUtils; // importacion de clase externa
import org.springframework.stereotype.Controller; // importacion de clase externa
import org.springframework.ui.Model; // importacion de clase externa
import org.springframework.validation.BindingResult; // importacion de clase externa
import org.springframework.web.bind.annotation.*; // importacion de clase externa
// importacion de clase externa
// importacion de clase externa
import org.springframework.web.servlet.mvc.support.RedirectAttributes;

@Controller // MVC (devuelve vistas)
@RequestMapping("/productos") // ruta base del controlador
@RequiredArgsConstructor
public class ProductoController { // clase publica

    private final ProductoService service; // inmutable

    /** GET /productos --> lista */
    @GetMapping // endpoint HTTP GET
    public String listar(Model model) { // metodo publico
        model.addAttribute("productos", service.listar());
        return "productos/lista"; // -> templates/productos/lista.html
    }

    /** GET /productos/nuevo --> formulario vacio */
    @GetMapping("/nuevo") // endpoint HTTP GET
    public String mostrarFormulario(Model model) { // metodo publico
        model.addAttribute("form", new ProductoForm());
        return "productos/formulario"; // retorna
    }

    /** POST /productos --> guardar */
    @PostMapping // endpoint HTTP POST (crear)
    // metodo publico
    // metodo publico
    public String guardar(@Valid @ModelAttribute("form") ProductoForm form,
                          BindingResult br,
                          RedirectAttributes ra) {
        if (br.hasErrors()) { // condicional
            // re-renderizar el formulario con los errores
            return "productos/formulario"; // retorna
        }
        Producto p = new Producto();
        BeanUtils.copyProperties(form, p);
        service.guardar(p);
        ra.addFlashAttribute("mensaje", "Producto guardado");
        return "redirect:/productos"; // patron PRG anti F5
    }

    /** GET /productos/{id}/editar */
    @GetMapping("/{id}/editar") // endpoint HTTP GET
    // metodo publico
    // metodo publico
    public String editar(@PathVariable Long id, Model model) {
        Producto p = service.buscar(id);
        ProductoForm form = new ProductoForm();
        BeanUtils.copyProperties(p, form);
        model.addAttribute("form", form);
        return "productos/formulario"; // retorna
    }

    /** POST /productos/{id}/eliminar */
    @PostMapping("/{id}/eliminar") // endpoint HTTP POST (crear)
    // metodo publico
    // metodo publico
    public String eliminar(@PathVariable Long id, RedirectAttributes ra) {
        service.eliminar(id);
        ra.addFlashAttribute("mensaje", "Producto eliminado");
        return "redirect:/productos"; // retorna
    }
}
\end{lstlisting}

\textbf{Vista Thymeleaf de listado (Listado~\ref{lst:14.7}):}
\begin{lstlisting}[style=html,caption={templates/productos/lista.html},label=lst:14.7]
<!DOCTYPE html>
<html xmlns:th="http://www.thymeleaf.org" lang="es-EC">
<head>
  <meta charset="UTF-8">
  <title>Productos UTEQ</title>
  <link rel="stylesheet" href="/css/app.css">
</head>
<body>
<header>
  <h1>Catalogo de productos</h1>
  <a th:href="@{/productos/nuevo}" class="btn">Nuevo producto</a>
</header>

<div th:if="${mensaje}" class="alerta-exito" th:text="${mensaje}"></div>

<table>
  <thead>
    <tr><th>ID</th><th>Nombre</th><th>Precio</th>
        <th>Stock</th><th>Disponible</th><th>Acciones</th></tr>
  </thead>
  <tbody>
    <tr th:each="p : ${productos}">
      <td th:text="${p.id}"></td>
      <td th:text="${p.nombre}"></td>
      <td th:text="${'$' + #numbers.formatDecimal(p.precio,1,2)}"></td>
      <td th:text="${p.stock}"></td>
      <td th:text="${p.estaDisponible() ? 'Si' : 'No'}"></td>
      <td>
        <a th:href="@{/productos/{id}/editar(id=${p.id})}">Editar</a>
        <form th:action="@{/productos/{id}/eliminar(id=${p.id})}"
              method="post" style="display:inline">
          <button type="submit"
                  onclick="return confirm('Confirmar eliminacion?')">
            Eliminar
          </button>
        </form>
      </td>
    </tr>
  </tbody>
</table>
</body>
</html>
\end{lstlisting}

\textbf{Vista del formulario (Listado~\ref{lst:14.8}):}
\begin{lstlisting}[style=html,caption={templates/productos/formulario.html},label=lst:14.8]
<!DOCTYPE html>
<html xmlns:th="http://www.thymeleaf.org" lang="es-EC">
<head><meta charset="UTF-8"><title>Producto</title></head>
<body>
<h1 th:text="${form.id == null ? 'Nuevo producto' : 'Editar producto'}">
</h1>

<form th:action="@{/productos}" th:object="${form}" method="post">
  <input type="hidden" th:field="*{id}">

  <div>
    <label>Nombre</label>
    <input type="text" th:field="*{nombre}" required>
    <span th:errors="*{nombre}" class="error"></span>
  </div>

  <div>
    <label>Precio</label>
    <input type="number" step="0.01" th:field="*{precio}" required>
    <span th:errors="*{precio}" class="error"></span>
  </div>

  <div>
    <label>Stock</label>
    <input type="number" th:field="*{stock}" min="0" required>
    <span th:errors="*{stock}" class="error"></span>
  </div>

  <button type="submit">Guardar</button>
  <a th:href="@{/productos}">Cancelar</a>
</form>
</body>
</html>
\end{lstlisting}

\textbf{Notas pedagógicas sobre este código:}
\begin{itemize}
\item \emph{Patrón PRG (Post-Redirect-Get):} tras un POST exitoso se
hace \texttt{redirect:/productos}, no se renderiza directamente la
vista. Esto evita el bug clásico de doble envío al pulsar F5.
\item \emph{Flash attributes:} \texttt{RedirectAttributes} permite
pasar mensajes entre el redirect y la vista sin contaminar la URL.
\item \emph{ProductoForm} en lugar de \texttt{Producto} en el
controller: separa la representación de la entidad y la del
formulario; permite validaciones específicas de UI.
\item \emph{El modelo no es anémico:} \texttt{estaDisponible()} es
una regla del dominio.
\item \emph{Thymeleaf-natural:} las plantillas son HTML válido que se
puede abrir en navegador sin servidor (modo preview).
\end{itemize}
\end{cajaejemplo}

\section{Variantes del MVC}

MVC puro no es la única opción: en la historia han surgido
variantes con ventajas específicas (MVP, MVVM, MVI). La tabla
siguiente las catalogará y dirá cuándo usar cada una.


\subsection{MVP (Model-View-Presenter)}

En MVP el Presenter reemplaza al Controller pero acopla más
fuertemente con la Vista. La Vista delega completamente la lógica de
presentación al Presenter. Común en aplicaciones Android pre-MVVM.

\subsection{MVVM (Model-View-ViewModel)}

En MVVM, el ViewModel es una abstracción de la vista que mantiene su
estado. La vista se enlaza al ViewModel por \emph{data binding}
bidireccional. Angular y Vue son herederos parciales de este enfoque.

\subsection{Flux y Redux}

Patrones unidireccionales para SPAs: Action $\to$ Dispatcher $\to$
Store $\to$ View. Resuelven el problema de complejidad del estado en
SPAs grandes (Facebook lo creó para resolver bugs en Messenger).

\subsection{Signals (Angular 17+, SolidJS)}

Modelo reactivo basado en \emph{signals}: variables observables que
notifican automáticamente cambios. Una evolución moderna del MVC
reactivo original de Reenskaug, ahora aplicada al frontend.

\section{Práctica integrada: CRUD MVC con debugging}

La práctica integrada de la semana 14 implementa el CRUD MVC
trazando explícitamente el flujo de cada petición, para que el
estudiante visualice cómo trabaja el framework por debajo.


\begin{cajaejercicio}{Ejercicio 14.1 --- CRUD MVC completo con trazado de la petición en IntelliJ}

\textbf{Objetivo:} construir un CRUD MVC server-side y aprender a
\emph{depurar} cada fase del ciclo de vida de la petición usando los
breakpoints y el inspector de IntelliJ IDEA Ultimate.

\subsubsection*{IDE: IntelliJ IDEA Ultimate}

IntelliJ IDEA Ultimate provee herramientas específicas para
trazar el flujo MVC (debugger con puntos de quiebre por capas,
consola Spring, herramienta HTTP Client). Esta práctica las usa
todas.


\subsubsection*{Paso 1 --- Crear el proyecto}

\texttt{File $\rightarrow$ New $\rightarrow$ Project $\rightarrow$
Spring Boot}. Configurar:
\begin{itemize}
\item Type: Maven, JDK: 21, Spring Boot 4.x.
\item Group: \texttt{ec.edu.uteq}; Artifact: \texttt{mvc-debug}.
\item Dependencies: \emph{Spring Web}, \emph{Thymeleaf}, \emph{Spring
Data JPA}, \emph{H2 Database}, \emph{Lombok}, \emph{Validation},
\emph{DevTools}.
\end{itemize}

\subsubsection*{Paso 2 --- Implementar el CRUD del Ejemplo 14.1}

Copiar la estructura completa del ejemplo 14.1 (entidad, DTO,
service, controller, dos vistas Thymeleaf).

Agregar en \texttt{application.yml} (Listado~\ref{lst:14.9}):
\begin{lstlisting}[style=yaml,numbers=none,caption={Ejemplo de YAML},label=lst:14.9]
spring: # configuracion de Spring
  datasource: # configuracion de origen de datos (BD)
    url: jdbc:h2:mem:appweb;DB_CLOSE_DELAY=-1 # URL JDBC de conexion a la BD
    driver-class-name: org.h2.Driver # clase del driver JDBC
    username: sa # usuario de la BD
    password: # contrasena (mejor leer de env var)
  jpa: # configuracion JPA/Hibernate
    hibernate.ddl-auto: create-drop # estrategia de generacion de schema
    show-sql: true # imprime SQL en consola (solo desarrollo)
    properties.hibernate.format_sql: true
  h2.console:
    enabled: true
    path:    /h2 # ruta donde se expone
\end{lstlisting}

\subsubsection*{Paso 3 --- Inicializar datos de prueba}

Crear \texttt{DatosIniciales.java} (Listado~\ref{lst:14.10}):
\begin{lstlisting}[style=java,numbers=none,caption={Ejemplo de Java},label=lst:14.10]
@Component // bean Spring generico
public class DatosIniciales implements CommandLineRunner { // clase publica
    private final ProductoRepository repo; // inmutable
    // metodo publico
    // metodo publico
    public DatosIniciales(ProductoRepository r) { this.repo = r; }
    @Override
    public void run(String... args) { // metodo publico
        repo.save(new Producto(null, "Cafe galletti", new BigDecimal("4.50"), 30));
        repo.save(new Producto(null, "Cacao Pacari",  new BigDecimal("6.00"),  5));
        repo.save(new Producto(null, "Te organico",   new BigDecimal("3.20"), 12));
    }
}
\end{lstlisting}

\subsubsection*{Paso 4 --- Ejecutar la aplicación}

Pulsar el botón verde de \emph{play} junto a la clase principal.
Esperar el log \texttt{Started McvDebugApplication}.

Abrir \url{http://localhost:8080/productos}: aparece el listado con
los tres productos iniciales.

\subsubsection*{Paso 5 --- Colocar 9 breakpoints estratégicos}

Para entender el ciclo de vida, colocar breakpoints (clic en el
margen izquierdo del editor) en:

\begin{enumerate}
\item \texttt{ProductoController.listar()} --- línea con
\texttt{model.addAttribute}.
\item \texttt{ProductoService.listar()} --- línea con
\texttt{return repo.findAll();}.
\item \texttt{ProductoController.mostrarFormulario()}.
\item \texttt{ProductoController.guardar()} --- línea con
\texttt{br.hasErrors()}.
\item \texttt{ProductoController.guardar()} --- línea con
\texttt{service.guardar(p);}.
\item \texttt{ProductoService.guardar(Producto)} --- la línea con
\texttt{repo.save(p);}.
\item \texttt{ProductoController.editar()}.
\item \texttt{ProductoController.eliminar()}.
\item \texttt{ProductoForm} --- la línea de la anotación
\texttt{@NotBlank} (para ver el orden de validación).
\end{enumerate}

\subsubsection*{Paso 6 --- Reiniciar en modo Debug}

Detener la app. Pulsar el ícono \emph{bug} (\texttt{Shift+F9}) para
arrancar en modo debug. La app está lista para depurar.

\subsubsection*{Paso 7 --- Trazar una petición GET (listar)}

Con la aplicación corriendo, el estudiante debe trazar una
petición \texttt{GET} completa observando cómo se mueve por cada
capa.


\begin{enumerate}
\item En el navegador, recargar
\url{http://localhost:8080/productos}.
\item IntelliJ se detiene en breakpoint 1 (\texttt{listar()}).
\item Observar el panel \emph{Variables}: aparece \texttt{model}
(vacío) y los parámetros del método.
\item Pulsar \texttt{F8} (Step Over): se ejecuta
\texttt{model.addAttribute(\dots, service.listar())}.
\item Pulsar \texttt{F7} (Step Into) en \texttt{service.listar()}:
salta al breakpoint 2.
\item En \emph{Variables} ver que \texttt{repo} es un proxy de Spring.
\item Pulsar \texttt{F8}: se ejecuta \texttt{repo.findAll()};
en la consola aparece el SQL: \texttt{select * from productos}.
\item Pulsar \texttt{F9} (Resume): la petición continúa, Thymeleaf
renderiza la vista, el navegador recibe el HTML.
\end{enumerate}

\subsubsection*{Paso 8 --- Trazar una petición POST con error de validación}

El segundo trazado es una petición \texttt{POST} con datos
inválidos: el objetivo es observar cómo \texttt{BindingResult}
captura los errores y la vista los muestra.


\begin{enumerate}
\item Ir a \url{http://localhost:8080/productos/nuevo}.
\item Llenar solo el campo «Stock» (dejar nombre y precio vacíos).
\item Pulsar Guardar.
\item IntelliJ se detiene en breakpoint 4 (\texttt{br.hasErrors()}).
\item Inspeccionar \texttt{br}: el panel \emph{Variables} muestra
una lista de errores (\texttt{nombre: NotBlank},
\texttt{precio: NotNull}).
\item Pulsar \texttt{F8}: como hay errores, retorna
\texttt{\char34{}productos/formulario\char34{}} con los errores marcados en la vista.
\end{enumerate}

\subsubsection*{Paso 9 --- Trazar un POST exitoso}

El tercer trazado es un POST exitoso: observa cómo se
guarda en BD y cómo se aplica el patrón POST-Redirect-GET para
evitar reenvíos del formulario.


\begin{enumerate}
\item Recargar el formulario; ahora llenar todos los campos
correctamente.
\item Pulsar Guardar.
\item IntelliJ se detiene en breakpoint 4 (\texttt{br.hasErrors()}).
\item \texttt{br.hasErrors()} es \texttt{false}; pulsar \texttt{F8}.
\item Salta al breakpoint 5: a punto de ejecutar
\texttt{service.guardar(p)}.
\item Pulsar \texttt{F7} (Step Into): salta al breakpoint 6
(\texttt{repo.save(p)}).
\item En \emph{Variables}, observar que \texttt{p.id} es \texttt{null}
ANTES del save.
\item Pulsar \texttt{F8}: \texttt{p.id} ahora tiene un valor (la BD
asignó la clave primaria).
\item En la consola, ver el \texttt{INSERT INTO productos} generado.
\item Pulsar \texttt{F9}: la respuesta es un redirect a
\texttt{/productos}, el navegador hace un GET nuevo, listamos
todos los productos (vuelve a parar en breakpoint 1).
\end{enumerate}

\subsubsection*{Paso 10 --- Inspeccionar la BD H2 en vivo}

La consola H2 permite inspeccionar la BD en tiempo real
durante la depuración. Esto es invaluable para verificar que los
datos llegaron a persistirse correctamente.


\begin{enumerate}
\item Sin detener la app, abrir \url{http://localhost:8080/h2} en otra
pestaña.
\item En \emph{JDBC URL} introducir: \texttt{jdbc:h2:mem:appweb}.
User: \texttt{sa}, sin password.
\item Pulsar Connect.
\item Ejecutar \texttt{SELECT * FROM PRODUCTOS;}: aparecen los
productos iniciales más los que se hayan creado.
\end{enumerate}

\subsubsection*{Resultado esperado}

El estudiante puede:
\begin{itemize}
\item Articular el flujo: \texttt{Controller $\to$ Service $\to$
Repository $\to$ BD $\to$ Service $\to$ Controller $\to$ Vista}.
\item Identificar dónde ocurre la validación
(\texttt{@Valid + BindingResult}).
\item Explicar el patrón PRG y por qué se usa.
\item Distinguir Modelo (Producto/Service/Repository) de Vista
(Thymeleaf) de Controller (handler HTTP).
\item Ver el SQL generado por JPA en tiempo real en la consola.
\end{itemize}

\subsubsection*{Reflexión final escrita}

Responder en máximo 250 palabras: ¿Qué pasaría si el controller
llamara directamente al \texttt{ProductoRepository} sin pasar por el
\texttt{ProductoService}? ¿Qué se rompería? ¿Qué clases serían más
difíciles de testear?

\textbf{Esta práctica es la \emph{Sumativa GA-Práctica en clase} de
la semana 14.}
\end{cajaejercicio}

\section{Buenas prácticas profesionales en MVC}

La caja siguiente recoge las reglas profesionales para
escribir MVC mantenible. Aplican tanto a Spring Boot como a
ASP.NET Core y Laravel.


\begin{cajabuenas}{Quince reglas profesionales 2026}
\begin{enumerate}
\item Controllers \emph{thin}: máximo 50 líneas por método de
controller, idealmente 10-20.
\item Toda la lógica de negocio en servicios, no en controllers ni
en vistas.
\item DTOs separados de entidades JPA: nunca exponer entidades
directamente al cliente.
\item Validación con Bean Validation (\texttt{@NotBlank},
\texttt{@Size}, \texttt{@Email}) sobre validación manual.
\item Patrón PRG (Post-Redirect-Get) para evitar reenvíos al pulsar
F5.
\item Flash attributes para mensajes de éxito/error en redirects.
\item Manejo global de excepciones con
\texttt{@ControllerAdvice}/\texttt{@RestControllerAdvice}.
\item Vistas (templates) sin lógica compleja: solo iteración,
condicionales simples y formato.
\item URLs RESTful: \texttt{/productos} para listar, \texttt{/productos/123}
para uno específico, métodos HTTP semánticos.
\item Repositorios definidos como interfaces (Spring Data) para
facilitar tests.
\item Servicios anotados \texttt{@Transactional} en los métodos que
mutan datos.
\item Logging estructurado en cada capa: controller (acceso),
servicio (negocio), repositorio (datos).
\item Tests unitarios de servicio con mock de repositorio; tests de
integración con MockMvc para controllers.
\item Documentación de la API REST con OpenAPI/Swagger.
\item Internacionalización (i18n) desde el inicio: archivos
\texttt{messages\_es.properties}, \texttt{messages\_en.properties}.
\end{enumerate}
\end{cajabuenas}

\section{Contraste bibliográfico de los conceptos clave}

La separación de responsabilidades que propone MVC se ha discutido
profundamente en la literatura. \cite{galloway2011mvc} es la
referencia clásica para entender MVC del lado servidor en
\texttt{ASP.NET}. \cite{walls2022spring} y \cite{spilca2024spring}
muestran cómo Spring MVC implementa el patrón con
\texttt{DispatcherServlet}, \texttt{HandlerMapping} y
\texttt{ViewResolver}. En el ámbito del frontend, las arquitecturas
modernas (Angular, React) han evolucionado hacia variantes
\emph{component-based} que rompen la estricta tripartición MVC pero
mantienen el principio de separación
\cite{richards2020fundamentals}. \cite{richardson2007restful} documenta
cómo el patrón \emph{API REST + SPA} es hoy el reemplazo dominante
del MVC server-side tradicional, con tres ventajas: equipos
desacoplados (frontend/backend), reutilización del backend para
varios clientes (web, móvil, IoT) y mejor experiencia interactiva.

\section{Ejercicio resuelto adicional con resultado visual}

Como cierre, el siguiente ejercicio resuelto adicional
implementa un CRUD completo con Thymeleaf y se complementa con un
propuesto que lo migra a SPA Angular.


\begin{cajaejemplo}{Ejercicio resuelto 14.1 --- App MVC de tareas con Spring Boot y Thymeleaf}
\textbf{Enunciado.} Implementar una aplicación MVC server-side de
gestión de tareas (\emph{to-do list}) con Spring Boot 4, Thymeleaf
y H2. Permitir listar, crear, completar y eliminar tareas.

\textbf{Modelo (\texttt{Tarea.java}):}
\begin{lstlisting}[language=Java,caption={Ejemplo de Java},label=lst:14.11]
@Entity @Table(name="tareas") // entidad JPA (mapea a tabla)
public class Tarea { // clase publica
  @Id @GeneratedValue(strategy=GenerationType.IDENTITY) // campo clave primaria
  private Long id; // campo privado
  @NotBlank @Size(max=120) private String titulo; // validacion: no nulo y no vacio
  private boolean completada; // campo privado
  // getters/setters
}
\end{lstlisting}

\textbf{Controlador (\texttt{TareaController.java}):}
\begin{lstlisting}[language=Java,caption={Ejemplo de Java},label=lst:14.12]
@Controller @RequestMapping("/tareas") // MVC (devuelve vistas)
public class TareaController { // clase publica
  private final TareaRepository repo; // inmutable
  // metodo publico
  // metodo publico
  public TareaController(TareaRepository r) { this.repo = r; }

  @GetMapping // endpoint HTTP GET
  public String listar(Model m) { // metodo publico
    m.addAttribute("tareas", repo.findAll());
    m.addAttribute("nueva", new Tarea());
    return "tareas"; // retorna
  }

  @PostMapping // endpoint HTTP POST (crear)
  // metodo publico
  // metodo publico
  public String crear(@Valid @ModelAttribute("nueva") Tarea t,
                      BindingResult br) {
    if (br.hasErrors()) return "tareas"; // condicional
    repo.save(t);
    return "redirect:/tareas"; // retorna
  }

  @PostMapping("/{id}/completar") // endpoint HTTP POST (crear)
  public String completar(@PathVariable Long id) { // metodo publico
    repo.findById(id).ifPresent(t -> {
      t.setCompletada(!t.isCompletada());
      repo.save(t);
    });
    return "redirect:/tareas"; // retorna
  }

  @PostMapping("/{id}/eliminar") // endpoint HTTP POST (crear)
  public String eliminar(@PathVariable Long id) { // metodo publico
    repo.deleteById(id);
    return "redirect:/tareas"; // retorna
  }
}
\end{lstlisting}

\textbf{Vista (\texttt{templates/tareas.html} --- Thymeleaf):}
\begin{lstlisting}[language=HTML5,caption={Ejemplo de HTML},label=lst:14.13]
<!DOCTYPE html>
<html xmlns:th="http://www.thymeleaf.org" lang="es">
<head><meta charset="UTF-8"><title>Mis tareas</title></head>
<body>
  <h1>Mis tareas</h1>
  <form th:action="@{/tareas}" th:object="${nueva}" method="post">
    <input type="text" th:field="*{titulo}" placeholder="Nueva tarea">
    <button type="submit">Agregar</button>
    <span th:errors="*{titulo}" style="color:red"></span>
  </form>
  <ul>
    <li th:each="t : ${tareas}">
      <span th:text="${t.titulo}"
            th:style="${t.completada} ? 'text-decoration:line-through' : ''"></span>
      <form th:action="@{|/tareas/${t.id}/completar|}"
            method="post" style="display:inline">
        <button type="submit"
                th:text="${t.completada} ? 'Reabrir' : 'Completar'"></button>
      </form>
      <form th:action="@{|/tareas/${t.id}/eliminar|}"
            method="post" style="display:inline">
        <button type="submit">Eliminar</button>
      </form>
    </li>
  </ul>
</body>
</html>
\end{lstlisting}

\textbf{Resultado visual} en el navegador:

\begin{center}
\fbox{\begin{minipage}{0.85\textwidth}
\vspace{0.2cm}
\hspace{0.3cm}\textbf{\Large Mis tareas}\\[0.4cm]
\hspace{0.3cm}\fbox{\texttt{Nueva tarea\ldots\ldots\ldots\ldots\ldots}} \fbox{Agregar}\\[0.5cm]
\hspace{0.3cm}$\bullet$ \emph{Leer cap.\ 14 del libro} \fbox{\small Reabrir} \fbox{\small Eliminar}\\[0.2cm]
\hspace{0.3cm}$\bullet$ Estudiar Spring MVC \fbox{\small Completar} \fbox{\small Eliminar}\\[0.2cm]
\hspace{0.3cm}$\bullet$ Hacer el ejercicio propuesto 14.2 \fbox{\small Completar} \fbox{\small Eliminar}\\[0.2cm]
\vspace{0.1cm}
\end{minipage}}
\end{center}

\textbf{Discusión.} El flujo MVC completo se ve en una sola pantalla:
el navegador hace \texttt{GET /tareas} $\to$
\texttt{DispatcherServlet} despacha a \texttt{TareaController.listar()}
$\to$ el controlador consulta el repositorio (modelo) y prepara los
datos en el \texttt{Model} $\to$ Thymeleaf renderiza la vista con esos
datos $\to$ el navegador recibe HTML. Al hacer clic en
\emph{Completar}, ocurre el patrón \emph{POST-Redirect-GET}
\cite{walls2022spring}: \texttt{POST} ejecuta la acción y el
servidor responde \texttt{302 Found} con \texttt{Location: /tareas},
forzando un nuevo \texttt{GET}. Esto evita el reenvío accidental del
formulario al recargar.
\end{cajaejemplo}

\begin{cajaejercicio}{Ejercicio propuesto 14.2 --- Migrar la app de tareas a SPA Angular + REST}
\textbf{Enunciado.} Tomar la aplicación MVC server-side del ejercicio
14.1 y refactorizarla a una arquitectura de dos niveles: backend Spring
Boot expone una API REST y un frontend Angular 19 SPA consume la API.
Comparar las dos arquitecturas en términos de líneas de código,
acoplamiento, experiencia de usuario y testabilidad.

\textbf{Requisitos:}
\begin{enumerate}
\item Convertir \texttt{TareaController} de \texttt{@Controller} a
\texttt{@RestController}, eliminando Thymeleaf.
\item Endpoints: \texttt{GET /api/v1/tareas},
\texttt{POST /api/v1/tareas}, \texttt{PATCH /api/v1/tareas/\{id\}/completar},
\texttt{DELETE /api/v1/tareas/\{id\}}.
\item Frontend Angular con un componente \texttt{TareasComponent}
que consuma la API mediante \texttt{HttpClient}.
\item Manejar CORS correctamente con
\texttt{@CrossOrigin(\char34{}http://localhost:4200\char34{})} o configuración
global.
\item Documentación OpenAPI con springdoc.
\item Test E2E con Playwright o Cypress que verifique el flujo
completo (crear, completar, eliminar).
\end{enumerate}

\textbf{Producto final:} un documento (1--2 páginas) que compare las
dos arquitecturas y argumente cuándo elegir cada una, citando al
menos a \cite{galloway2011mvc} y \cite{richardson2007restful}.
\end{cajaejercicio}

% =====================================================================
%       SEMANA 15: SERVICIOS WEB REST Y SOAP
% =====================================================================
\chapter{Servicios web: arquitecturas REST, SOAP y diseño de APIs}
\label{cap:semana15}

\epigraph{«Una buena API es como una buena novela: invita a explorarla,
revela su lógica gradualmente y nunca traiciona al lector con
inconsistencias arbitrarias.»}{--- \textit{Adaptado de Joshua Bloch}}

\section*{Resultado de aprendizaje de la semana}
Al concluir la semana 15, el estudiante diseña, implementa, documenta
y consume servicios web aplicando el estilo arquitectónico REST según
Roy Fielding \cite{fielding2000rest}, distingue REST de SOAP en sus
casos de uso reales, integra autenticación con JWT y produce
documentación OpenAPI 3.0 conforme a estándares profesionales de la
industria 2026.

\section{Historia de los servicios web}

Los servicios web son la forma estándar de comunicación
entre sistemas distintos. Su historia atraviesa tres generaciones:
SOAP (1998), REST (2000) y GraphQL/gRPC (2015). Conocerlas permite
elegir adecuadamente para cada proyecto.


\begin{cajahistoria}{De RPC a REST --- treinta años de evolución de las APIs distribuidas}
La idea de que un programa pueda invocar funciones de otro a través
de la red precede a la web. En 1981, Bruce Nelson y Andrew Birrell
formalizaron el \emph{Remote Procedure Call} (RPC) en el seminario
del que surgió la primera implementación seria, Sun RPC. El sueño:
hacer que llamar a una función remota sea sintácticamente idéntico
a llamar a una función local.

A finales de los noventa, ese sueño se reformuló en la web. Microsoft
junto con DevelopMentor publicaron en 1998 el primer borrador de
\textbf{SOAP} (Simple Object Access Protocol): RPC sobre HTTP usando
XML como formato de mensaje. SOAP se complementó con \textbf{WSDL}
(Web Services Description Language) en 2001 y con \textbf{UDDI} para
descubrimiento de servicios. La trinidad SOAP+WSDL+UDDI prometía un
ecosistema empresarial donde cualquier servicio sería descubrible y
consumible automáticamente.

La realidad fue más áspera. SOAP era verboso (envoltorios XML
enormes), el ecosistema de \emph{WS-*} (WS-Security, WS-ReliableMessaging,
WS-Transaction\dots) creció hasta volverse impracticable, y los
clientes JavaScript del navegador no podían consumirlo elegantemente.

En 2000 \textbf{Roy Fielding} —uno de los autores principales de
HTTP/1.1 y cofundador de Apache— defendió en su tesis doctoral el
estilo arquitectónico \textbf{REST} (Representational State Transfer)
\cite{fielding2000rest}. REST no era un protocolo: era un conjunto de
\emph{restricciones arquitectónicas} (cliente-servidor, stateless,
cacheable, interfaz uniforme, sistema en capas, código bajo demanda
opcional) que aprovechaban HTTP en lugar de pelearse con él. La
trayectoria a partir de 2005 fue rápida: las APIs públicas masivas
(Flickr, Twitter, Amazon S3, GitHub) eligieron REST. Para 2015, REST
era el estilo dominante en la web. SOAP quedó relegado a sistemas
empresariales legados y a sectores conservadores como banca y salud.

En 2026 conviven cuatro paradigmas de servicios web: \textbf{REST}
(dominante en APIs públicas y SPAs), \textbf{GraphQL} (creado por
Facebook en 2012, popular en frontends complejos), \textbf{gRPC}
(Google, 2015, eficiente en comunicación servicio-servicio) y
\textbf{SOAP} (legado pero vigente en banca, salud, gobierno).
\end{cajahistoria}

\subsection{Tabla evolutiva de los servicios web}

La tabla siguiente recoge los hitos cronológicos de los
servicios web, desde XML-RPC hasta GraphQL y gRPC.


\begin{table}[htbp]
\centering
\caption{Hitos en la historia de los servicios web.}
\footnotesize
\begin{tabularx}{\textwidth}{lll}
\toprule
\textbf{Año} & \textbf{Hito} & \textbf{Aporte}\\
\midrule
1981 & Sun RPC & Primer RPC industrial.\\
1991 & CORBA & Estándar OMG; complejo, propietario.\\
1998 & SOAP 1.0 & RPC sobre HTTP+XML; Microsoft, IBM.\\
2000 & Tesis de Fielding & REST formalizado como estilo arquitectónico.\\
2001 & WSDL 1.1 & Contrato XML para SOAP.\\
2003 & SOAP 1.2 (W3C) & Recomendación oficial de SOAP.\\
2005 & Atom Publishing Protocol & Primer estándar REST formal.\\
2008 & RESTful API patterns & Modelo de Madurez de Richardson.\\
2009 & Twitter API REST & APIs públicas masivas adoptan REST.\\
2012 & GraphQL (Facebook) & Cliente decide qué campos pedir.\\
2014 & OpenAPI 2.0 (Swagger) & Estándar para describir APIs REST.\\
2015 & gRPC (Google) & RPC moderno con HTTP/2 y Protocol Buffers.\\
2017 & OpenAPI 3.0 & Versión moderna del estándar.\\
2021 & OpenAPI 3.1 & Compatibilidad total con JSON Schema 2020-12.\\
2024 & gRPC-Web estable & gRPC accesible desde navegador.\\
\bottomrule
\end{tabularx}
\end{table}

\section{Las seis restricciones de la arquitectura REST}

Roy Fielding identificó seis restricciones que un sistema debe
cumplir para considerarse RESTful. No son sugerencias: son
restricciones arquitectónicas con consecuencias específicas.

\begin{cajadefinicion}{Las restricciones REST de Fielding (2000)}
\begin{enumerate}
\item \textbf{Cliente-servidor.} Separación clara de
responsabilidades: el cliente no conoce la persistencia, el servidor
no conoce la presentación.

\item \textbf{Stateless.} El servidor NO guarda estado de sesión
entre peticiones. Cada petición lleva toda la información necesaria
para procesarla. Permite escalar horizontalmente sin sticky sessions.

\item \textbf{Cacheable.} Las respuestas indican explícitamente si
pueden cachearse y por cuánto tiempo (\texttt{Cache-Control},
\texttt{ETag}). Mejora rendimiento y escalabilidad.

\item \textbf{Interfaz uniforme.} Cuatro sub-restricciones:
\begin{itemize}
\item Identificación de recursos por URI.
\item Manipulación a través de representaciones (JSON, XML).
\item Mensajes auto-descriptivos (cabeceras explícitas).
\item HATEOAS (Hypermedia as the Engine of Application State):
los recursos enlazan a otros relacionados.
\end{itemize}

\item \textbf{Sistema en capas.} El cliente no sabe si habla con el
servidor final o con un proxy/balanceador. Permite intermediarios
(CDN, gateways, balanceadores).

\item \textbf{Código bajo demanda} (opcional). El servidor puede
extender al cliente enviando código ejecutable (típicamente
JavaScript).
\end{enumerate}
\end{cajadefinicion}

\subsection{Modelo de Madurez de Richardson}

Leonard Richardson propuso en 2008 un modelo escalonado para evaluar
cuán RESTful es una API:

\begin{table}[htbp]
\centering
\caption{Modelo de Madurez de Richardson para APIs REST.}
\footnotesize
\begin{tabularx}{\textwidth}{llX}
\toprule
\textbf{Nivel} & \textbf{Nombre} & \textbf{Características}\\
\midrule
0 & «The Swamp of POX» & XML/JSON sobre HTTP, un solo endpoint, todo es POST.\\
1 & Recursos & URLs distintas por recurso (\texttt{/usuarios/42}), pero todo POST.\\
2 & Verbos HTTP & Uso semántico de GET, POST, PUT, DELETE; códigos correctos.\\
3 & HATEOAS & Las respuestas incluyen enlaces de navegación a recursos relacionados.\\
\bottomrule
\end{tabularx}
\end{table}

\textbf{En 2026, la mayoría de APIs profesionales se sitúan en
Nivel 2}. El nivel 3 (HATEOAS) es elegante en teoría pero
infrecuente en la práctica: añade complejidad sin que la mayoría de
clientes lo aproveche realmente.

\section{Diseño de URIs y verbos HTTP}

La definición de las URIs y el uso de los verbos HTTP son
las dos decisiones más visibles de una API REST. Hacerlas mal se
nota desde el primer endpoint.


\subsection{Convenciones de nomenclatura}

Las URIs de una API REST siguen convenciones precisas: son
sustantivos en plural, sin verbos, jerárquicas. La lista siguiente
las recoge con ejemplos.

\textbf{Ejemplo corto comentado (URIs REST canónicas) (Listado~\ref{lst:15.1}):}
\begin{lstlisting}[language=bash,numbers=none,caption={Ejemplo de Shell},label=lst:15.1]
# BIEN: sustantivo en plural, jerarquia clara
GET    /api/v1/usuarios # listar todos
GET    /api/v1/usuarios/42 # uno especifico
POST   /api/v1/usuarios # crear nuevo
PUT    /api/v1/usuarios/42 # actualizar completo
PATCH  /api/v1/usuarios/42 # actualizar parcial
DELETE /api/v1/usuarios/42 # eliminar
GET    /api/v1/usuarios/42/pedidos # recurso anidado

# MAL: verbos en URI, singular, sin version
GET    /getUsuario?id=42 # verbo en URI
POST   /crearUsuario # accion como recurso
GET    /usuario/42 # singular en lugar de plural
\end{lstlisting}


\begin{cajabuenas}{Doce reglas para URIs RESTful profesionales}
\begin{enumerate}
\item \textbf{Sustantivos en plural} para colecciones:
\texttt{/productos}, no \texttt{/getProducto}.
\item \textbf{Recursos jerárquicos}:
\texttt{/usuarios/42/pedidos/7/items}.
\item \textbf{Sin verbos en la URI}: el verbo va en el método HTTP.
\item \textbf{Minúsculas siempre}, palabras separadas por guiones:
\texttt{/ordenes-pendientes}, no \texttt{/OrdenesPendientes}.
\item \textbf{Sin extensiones de archivo}: \texttt{/usuarios/42}, no
\texttt{/usuarios/42.json}.
\item \textbf{Versionado en la URL} con prefijo \texttt{/api/v1/}:
\texttt{/api/v1/productos}.
\item \textbf{Filtros por query string}:
\texttt{/productos?categoria=cafe\&min\_precio=2}.
\item \textbf{Paginación estándar}:
\texttt{?page=3\&size=20} o \texttt{?offset=60\&limit=20}.
\item \textbf{Ordenamiento por query}:
\texttt{?sort=precio,-nombre} (\texttt{-} para descendente).
\item \textbf{Búsqueda con parámetro \texttt{q}}:
\texttt{?q=galletti}.
\item \textbf{IDs estables} (UUIDs preferidos sobre auto-incrementales
en APIs públicas).
\item \textbf{Sub-recursos para acciones}: \texttt{/pedidos/7/cancelar}
con POST.
\end{enumerate}
\end{cajabuenas}

\subsection{Mapeo verbo-acción canónico}

Cada verbo HTTP tiene un significado canónico y debe
mapearse a la acción CRUD correspondiente. Confundirlos rompe la
semántica REST y genera APIs confusas.


\begin{table}[htbp]
\centering
\caption{Mapeo verbos HTTP a operaciones CRUD.}
\footnotesize
\begin{tabularx}{\textwidth}{llXX}
\toprule
\textbf{Verbo} & \textbf{URI} & \textbf{Operación} & \textbf{Códigos típicos}\\
\midrule
GET & /productos & Listar & 200 OK \\
GET & /productos/42 & Obtener uno & 200, 404 \\
POST & /productos & Crear & 201 Created (con \texttt{Location}) \\
PUT & /productos/42 & Reemplazar completo & 200, 204 \\
PATCH & /productos/42 & Actualización parcial & 200, 204 \\
DELETE & /productos/42 & Eliminar & 204 No Content, 404 \\
HEAD & /productos/42 & Solo cabeceras & 200, 404 \\
OPTIONS & /productos & Verbos soportados & 200 (CORS) \\
\bottomrule
\end{tabularx}
\end{table}

\subsection{Idempotencia y seguridad}

Dos propiedades cruciales de los métodos HTTP:

\begin{itemize}
\item \textbf{Seguro} (\emph{safe}): no modifica estado del servidor.
GET, HEAD, OPTIONS son seguros.
\item \textbf{Idempotente}: ejecutar la operación N veces produce el
mismo efecto que ejecutarla una vez. GET, PUT, DELETE, HEAD, OPTIONS
son idempotentes. POST y PATCH \emph{no} lo son.
\end{itemize}

Los \emph{proxies} y reintentos automáticos asumen estas propiedades:
si su API rompe estas reglas, los clientes sufrirán bugs sutiles.

\section{Implementación de APIs REST en las tres plataformas}

Las tres plataformas del curso (Spring Boot, ASP.NET Core,
PHP) construyen APIs REST con APIs limpias y modernas. Las
subsecciones siguientes muestran la implementación canónica en cada
una.


\subsection{Spring Boot 4 con Spring Web}

Spring Boot expone APIs REST con \texttt{@RestController} y
anotaciones por método HTTP. Es la forma idiomática en el ecosistema
Java.


\begin{cajaejemplo}{Ejemplo 15.1 --- API REST de pedidos completa en Spring Boot}

\textbf{Enunciado.} Implementar una API REST de productos con Spring Boot, autenticación JWT, documentación OpenAPI 3.0 y Swagger UI interactivo.

\begin{lstlisting}[style=java,caption={PedidoController.java},label=lst:15.2]
package ec.edu.uteq.appweb.controller; // paquete del archivo

import ec.edu.uteq.appweb.dto.*; // importacion de clase externa
import ec.edu.uteq.appweb.service.PedidoService; // importacion de clase externa
import io.swagger.v3.oas.annotations.Operation; // importacion de clase externa
import io.swagger.v3.oas.annotations.tags.Tag; // importacion de clase externa
import jakarta.validation.Valid; // importacion de clase externa
import lombok.RequiredArgsConstructor; // importacion de clase externa
import org.springframework.data.domain.Page; // importacion de clase externa
import org.springframework.data.domain.Pageable; // importacion de clase externa
import org.springframework.http.ResponseEntity; // importacion de clase externa
import org.springframework.web.bind.annotation.*; // importacion de clase externa
import java.net.URI; // importacion de clase externa

@RestController // expone endpoints REST con JSON
@RequestMapping("/api/v1/pedidos") // ruta base del controlador
@RequiredArgsConstructor
// agrupa endpoints en Swagger UI
// agrupa endpoints en Swagger UI
@Tag(name = "Pedidos", description = "Gestion de pedidos del sistema")
public class PedidoController { // clase publica

    private final PedidoService service; // inmutable

    @GetMapping // endpoint HTTP GET
    @Operation(summary = "Lista paginada de pedidos") // OpenAPI
    public Page<PedidoDto> listar( // metodo publico
            @RequestParam(required = false) String estado, // query param
            Pageable pageable) {
        return service.buscar(estado, pageable); // retorna
    }

    @GetMapping("/{id}") // endpoint HTTP GET
    @Operation(summary = "Obtener pedido por ID") // OpenAPI
    // metodo publico
    // metodo publico
    public ResponseEntity<PedidoDto> obtener(@PathVariable Long id) {
        return ResponseEntity.ok(service.obtener(id)); // retorna
    }

    @PostMapping // endpoint HTTP POST (crear)
    @Operation(summary = "Crear nuevo pedido") // OpenAPI
    public ResponseEntity<PedidoDto> crear( // metodo publico
            @Valid @RequestBody PedidoCreateRequest req) { // activa validacion Bean Validation
        PedidoDto creado = service.crear(req);
        return ResponseEntity // retorna
            .created(URI.create("/api/v1/pedidos/" + creado.id()))
            .body(creado);
    }

    @PatchMapping("/{id}") // endpoint HTTP PATCH (actualizar parcial)
    @Operation(summary = "Actualizacion parcial") // OpenAPI
    public PedidoDto actualizar(@PathVariable Long id, // metodo publico
                                // activa validacion Bean Validation
                                // activa validacion Bean Validation
                                @Valid @RequestBody PedidoUpdateRequest req) {
        return service.actualizar(id, req); // retorna
    }

    @PostMapping("/{id}/cancelar") // endpoint HTTP POST (crear)
    @Operation(summary = "Cancelar un pedido") // OpenAPI
    public PedidoDto cancelar(@PathVariable Long id, // metodo publico
                              // query param
                              // query param
                              @RequestParam(required = false) String motivo) {
        return service.cancelar(id, motivo); // retorna
    }

    @DeleteMapping("/{id}") // endpoint HTTP DELETE (eliminar)
    @Operation(summary = "Eliminar pedido") // OpenAPI
    // metodo publico
    // metodo publico
    public ResponseEntity<Void> eliminar(@PathVariable Long id) {
        service.eliminar(id);
        return ResponseEntity.noContent().build(); // retorna
    }
}
\end{lstlisting}

\textbf{Manejo global de excepciones (RFC 7807) (Listado~\ref{lst:15.3}):}
\begin{lstlisting}[style=java,caption={GlobalExceptionHandler.java},label=lst:15.3]
package ec.edu.uteq.appweb.exception; // paquete del archivo

import jakarta.servlet.http.HttpServletRequest; // importacion de clase externa
import org.springframework.http.HttpStatus; // importacion de clase externa
import org.springframework.http.ProblemDetail; // importacion de clase externa
import org.springframework.http.ResponseEntity; // importacion de clase externa
// importacion de clase externa
// importacion de clase externa
import org.springframework.web.bind.MethodArgumentNotValidException;
import org.springframework.web.bind.annotation.*; // importacion de clase externa
import java.net.URI; // importacion de clase externa
import java.time.Instant; // importacion de clase externa
import java.util.*; // importacion de clase externa

@RestControllerAdvice // expone endpoints REST con JSON
public class GlobalExceptionHandler { // clase publica

    @ExceptionHandler(NoSuchElementException.class)
    public ResponseEntity<ProblemDetail> noEncontrado( // metodo publico
            NoSuchElementException ex, HttpServletRequest req) {
        ProblemDetail pd = ProblemDetail.forStatusAndDetail(
            HttpStatus.NOT_FOUND, ex.getMessage());
        pd.setType(URI.create("https://uteq.edu.ec/errors/no-encontrado"));
        pd.setTitle("Recurso no encontrado");
        pd.setInstance(URI.create(req.getRequestURI()));
        pd.setProperty("timestamp", Instant.now());
        // retorna
        // retorna
        return ResponseEntity.status(HttpStatus.NOT_FOUND).body(pd);
    }

    @ExceptionHandler(MethodArgumentNotValidException.class)
    public ResponseEntity<ProblemDetail> validacion( // metodo publico
            MethodArgumentNotValidException ex, HttpServletRequest req) {
        Map<String, String> errores = new LinkedHashMap<>();
        ex.getBindingResult().getFieldErrors().forEach(e ->
            errores.put(e.getField(), e.getDefaultMessage()));

        ProblemDetail pd = ProblemDetail.forStatusAndDetail(
            HttpStatus.UNPROCESSABLE_ENTITY,
            "Errores de validacion en " + errores.size() + " campo(s)");
        pd.setType(URI.create("https://uteq.edu.ec/errors/validacion"));
        pd.setTitle("Datos invalidos");
        pd.setInstance(URI.create(req.getRequestURI()));
        pd.setProperty("errores", errores);
        // retorna
        // retorna
        return ResponseEntity.unprocessableEntity().body(pd);
    }
}
\end{lstlisting}
\end{cajaejemplo}

\subsection{ASP.NET Core 10}

ASP.NET Core utiliza el atributo \texttt{[ApiController]} y
la herencia de \texttt{ControllerBase} para construir APIs REST.
Combinado con \emph{model binding} y \emph{model validation}
automáticos, ofrece una experiencia muy productiva.


\begin{cajaejemplo}{Ejemplo 15.2 --- Mismo recurso en ASP.NET Core 10}

\textbf{Enunciado.} Implementar el mismo recurso de productos visto antes pero en ASP.NET Core 10, para comparar idiomas: atributos \texttt{[ApiController]}, \texttt{[HttpGet]}, retornos \texttt{Ok()}/\texttt{NotFound()}.

\begin{lstlisting}[style=cs,caption={PedidosController.cs},label=lst:15.4]
using Microsoft.AspNetCore.Mvc; // importa namespace
using AppwebApi.Services; // importa namespace
using AppwebApi.Dtos; // importa namespace

namespace AppwebApi.Controllers; // namespace del archivo

[ApiController] // controlador de Web API (validacion automatica)
[Route("api/v1/[controller]")] // ruta base del controlador
[Produces("application/json")]
public class PedidosController : ControllerBase // clase publica
{
    private readonly IPedidoService _svc; // metodo o miembro privado

    // metodo o miembro publico
    // metodo o miembro publico
    public PedidosController(IPedidoService svc) => _svc = svc;

    /// <summary>Lista paginada de pedidos</summary>
    [HttpGet] // endpoint HTTP GET
    [ProducesResponseType(typeof(PageDto<PedidoDto>), 200)]
    // metodo o miembro publico
    // metodo o miembro publico
    public async Task<ActionResult<PageDto<PedidoDto>>> Listar(
        // mapea query string
        // mapea query string
        [FromQuery] string? estado, [FromQuery] int page = 1,
        [FromQuery] int size = 20) // mapea query string
        => Ok(await _svc.BuscarAsync(estado, page, size));

    [HttpGet("{id:long}")] // endpoint HTTP GET
    [ProducesResponseType(typeof(PedidoDto), 200)]
    [ProducesResponseType(404)]
    // metodo o miembro publico
    // metodo o miembro publico
    public async Task<ActionResult<PedidoDto>> Obtener(long id)
    {
        var p = await _svc.ObtenerAsync(id);
        return p is null ? NotFound() : Ok(p); // retorna
    }

    [HttpPost] // endpoint HTTP POST (crear)
    [ProducesResponseType(typeof(PedidoDto), 201)]
    [ProducesResponseType(422)]
    public async Task<ActionResult<PedidoDto>> Crear( // metodo o miembro publico
        [FromBody] PedidoCreateRequest req) // mapea body JSON al parametro
    {
        var creado = await _svc.CrearAsync(req);
        return CreatedAtAction(nameof(Obtener), // HTTP 201 (recurso creado)
            new { id = creado.Id }, creado);
    }

    [HttpDelete("{id:long}")] // endpoint HTTP DELETE (eliminar)
    [ProducesResponseType(204)]
    [ProducesResponseType(404)]
    public async Task<IActionResult> Eliminar(long id) // metodo o miembro publico
    {
        await _svc.EliminarAsync(id);
        return NoContent(); // HTTP 204 (sin cuerpo)
    }
}
\end{lstlisting}
\end{cajaejemplo}

\subsection{PHP 8.4 con Slim Framework 4}

PHP construye APIs REST profesionales mediante microframeworks
como Slim Framework 4. El siguiente ejemplo construye una API
minimalista pero idiomática.


\begin{cajaejemplo}{Ejemplo 15.3 --- API REST minimalista en PHP con Slim 4}

\textbf{Enunciado.} Construir una API REST minimalista en PHP usando el microframework Slim 4, ideal para servicios pequeños o microservicios PHP.

\begin{lstlisting}[style=php,caption={public/index.php},label=lst:15.5]
<?php
declare(strict_types=1);
require __DIR__ . '/../vendor/autoload.php'; // incluye otro archivo PHP (fatal si falla)

use Slim\Factory\AppFactory; // importa clase del namespace
use Psr\Http\Message\ServerRequestInterface as Req; // importa clase del namespace
use Psr\Http\Message\ResponseInterface as Res; // importa clase del namespace

$app = AppFactory::create();
$app->addBodyParsingMiddleware();

// GET /api/v1/pedidos
$app->get('/api/v1/pedidos', function (Req $req, Res $res) {
    $pdo = obtenerPDO();
    $estado = $req->getQueryParams()['estado'] ?? null;
    $sql = "SELECT id, total, estado, creado FROM pedidos";
    $params = [];
    // condicional
    // condicional
    if ($estado) { $sql .= " WHERE estado = ?"; $params[] = $estado; }
    $sql .= " ORDER BY creado DESC LIMIT 50";
    $stmt = $pdo->prepare($sql); // prepara consulta SQL (defensa anti-SQLi)
    $stmt->execute($params); // ejecuta la consulta con parametros enlazados
    // recupera todas las filas restantes
    // recupera todas las filas restantes
    $res->getBody()->write(json_encode(['data' => $stmt->fetchAll()]));
    // retorna valor
    // retorna valor
    return $res->withHeader('Content-Type','application/json');
});

// POST /api/v1/pedidos
$app->post('/api/v1/pedidos', function (Req $req, Res $res) {
    $body = $req->getParsedBody() ?? [];
    // condicional
    // condicional
    if (empty($body['total']) || !is_numeric($body['total'])) {
        $err = ['type'=>'/errors/validacion',
                'title'=>'Datos invalidos','status'=>422,
                'errors'=>['total'=>'es requerido y numerico']];
        $res->getBody()->write(json_encode($err));
        return $res->withStatus(422) // retorna valor
            ->withHeader('Content-Type','application/problem+json');
    }
    $pdo = obtenerPDO();
    $stmt = $pdo->prepare( // prepara consulta SQL (defensa anti-SQLi)
        "INSERT INTO pedidos (total, estado, creado) VALUES (?,?,NOW())");
    $stmt->execute([$body['total'], 'PENDIENTE']); // ejecuta la consulta con parametros enlazados
    $id = $pdo->lastInsertId();
    $data = ['id'=>$id,'total'=>$body['total'],'estado'=>'PENDIENTE'];
    $res->getBody()->write(json_encode($data));
    return $res->withStatus(201) // retorna valor
        ->withHeader('Content-Type','application/json')
        ->withHeader('Location',"/api/v1/pedidos/$id");
});

$app->run();
\end{lstlisting}
\end{cajaejemplo}

\section{REST vs SOAP: la comparativa definitiva}

REST y SOAP son los dos modelos arquitectónicos de servicios
web más utilizados. Comparar sus características objetivamente es
clave para elegir correctamente en cada proyecto.


\begin{cajacompara}{REST vs SOAP --- comparación práctica}
\centering\footnotesize
\begin{tabularx}{\textwidth}{lXX}
\toprule
\textbf{Criterio} & \textbf{REST (mainstream 2026)} & \textbf{SOAP (legado)}\\
\midrule
Formato & JSON (texto compacto) & XML envelope con header y body (verboso)\\
Transporte & HTTP/HTTPS & HTTP, SMTP, FTP, JMS\\
Contrato & OpenAPI 3.0 (opcional) & WSDL (obligatorio)\\
Estado & Stateless por restricción & Stateful o stateless\\
Tamaño mensaje & 3--5x más compacto & Pesado\\
Seguridad & HTTPS + JWT + OAuth 2.0 & WS-Security (más complejo)\\
Desde JavaScript & Trivial (fetch nativo) & Difícil (requiere SOAP client)\\
Caché HTTP nativa & Sí & No (POSTs)\\
Casos típicos & APIs públicas, microservicios, SPA & Banca, salud, B2B legado\\
Curva de aprendizaje & Baja & Alta\\
Tooling 2026 & Excelente & Maduro pero envejecido\\
\bottomrule
\end{tabularx}
\end{cajacompara}

\subsection{Cuándo aún tiene sentido SOAP en 2026}

Aunque REST es la opción dominante en 2026, hay escenarios
donde SOAP todavía tiene ventajas reales. La lista siguiente los
recoge.


\begin{itemize}
\item \textbf{Sistemas financieros nacionales}: muchos sistemas
bancarios mantienen interfaces SOAP por inercia regulatoria. En
Ecuador, los servicios SPI del Banco Central usan SOAP.
\item \textbf{Salud}: HL7 v3 sobre SOAP sigue vigente en hospitales.
\item \textbf{B2B con grandes corporaciones}: SAP, Oracle EBS y
sistemas similares exponen SOAP por defecto.
\item \textbf{Cuando se necesita WS-AtomicTransaction} (transacciones
distribuidas estrictas).
\end{itemize}

\subsection{Ejemplo SOAP mínimo}

Aunque el curso no usa SOAP en el PFC, conviene haber visto
al menos una petición SOAP completa para reconocerla cuando aparezca
en código heredado.

\textbf{Ejemplo corto comentado (petición SOAP) (Listado~\ref{lst:15.6}):}
\begin{lstlisting}[language=HTML5,numbers=none,caption={Ejemplo de HTML},label=lst:15.6]
<!-- POST /servicios/clima HTTP/1.1                              -->
<!-- Content-Type: text/xml; charset=utf-8                       -->
<!-- SOAPAction: "obtenerTemperatura"                            -->
<?xml version="1.0"?>
<!-- Envelope: contenedor obligatorio de todo SOAP                -->
<soap:Envelope xmlns:soap="http://schemas.xmlsoap.org/soap/envelope/">
  <!-- Body: contenido principal de la peticion                   -->
  <soap:Body>
    <!-- Operacion definida en el WSDL                            -->
    <obtenerTemperatura xmlns="http://uteq.edu.ec/clima">
      <ciudad>Quevedo</ciudad>
    </obtenerTemperatura>
  </soap:Body>
</soap:Envelope>
\end{lstlisting}


\begin{cajaejemplo}{Ejemplo 15.4 --- Cliente SOAP en PHP que consume servicio del SRI (estructura conceptual)}

\textbf{Enunciado.} Mostrar la estructura de un cliente SOAP en PHP que consume un servicio del SRI ecuatoriano (caso real del contexto local), demostrando el protocolo en producción.

\begin{lstlisting}[style=php,caption={consultar-ruc.php},label=lst:15.7]
<?php
declare(strict_types=1);

// PHP tiene cliente SOAP nativo: SoapClient
$wsdl = 'https://servicios.sri.gob.ec/aut.../wsdl';

try {
    $client = new SoapClient($wsdl, [
        'soap_version' => SOAP_1_2,
        'trace'        => true,
        'exceptions'   => true,
        'connection_timeout' => 10,
    ]);

    // Llamada SOAP: el metodo se invoca como si fuera local
    $resp = $client->consultarRuc(['numeroRuc' => '0992345678001']);

    print_r($resp);

    // Para debug del XML enviado/recibido:
    // echo $client->__getLastRequest();
    // echo $client->__getLastResponse();

} catch (SoapFault $e) {
    error_log("SOAP fault: " . $e->getMessage());
    http_response_code(502);
    // imprime al output del servidor
    // imprime al output del servidor
    echo json_encode(['error' => 'Servicio externo no disponible']);
}
\end{lstlisting}
\end{cajaejemplo}

\section{Documentación con OpenAPI 3.0 (Swagger)}

Una API sin documentación es inutilizable; una con
documentación a mano queda desincronizada al primer cambio.
OpenAPI 3.0 + Swagger UI resuelve esto generando documentación viva
desde el propio código.


\subsection{¿Qué es OpenAPI?}

OpenAPI 3.0 (antes Swagger) es la especificación estándar de la
industria para describir APIs REST. Permite:

\begin{itemize}
\item Documentación interactiva con Swagger UI.
\item Generación automática de clientes en cualquier lenguaje.
\item Pruebas y mocks automáticos.
\item Validación de contratos (contract testing).
\end{itemize}

\subsection{Configurar OpenAPI en Spring Boot}

Springdoc es la biblioteca que añade OpenAPI a Spring Boot
con configuración mínima. El siguiente fragmento la habilita y
expone Swagger UI.


\begin{cajaejemplo}{Ejemplo 15.5 --- OpenAPI con springdoc}

\textbf{Enunciado.} Configurar springdoc-openapi en una aplicación Spring Boot para exponer documentación OpenAPI 3.0 automática y la interfaz Swagger UI interactiva.

\textbf{En \texttt{pom.xml}:}
\begin{lstlisting}[style=xml,numbers=none,caption={Ejemplo de XML},label=lst:15.8]
<dependency>
  <groupId>org.springdoc</groupId>
  <artifactId>springdoc-openapi-starter-webmvc-ui</artifactId>
  <version>2.6.0</version>
</dependency>
\end{lstlisting}

\textbf{Configuración \texttt{OpenApiConfig.java}:}
\begin{lstlisting}[style=java,numbers=none,caption={Ejemplo de Java},label=lst:15.9]
@Configuration // clase de configuracion Spring
public class OpenApiConfig { // clase publica
    @Bean
    public OpenAPI api() { // metodo publico
        return new OpenAPI() // retorna
            .info(new Info()
                .title("API Aplicaciones Web UTEQ")
                .version("1.0.0")
                .description("API REST del PFC")
                .contact(new Contact().email("ti@uteq.edu.ec")))
            .components(new Components()
                .addSecuritySchemes("bearer-jwt",
                    new SecurityScheme()
                        .type(SecurityScheme.Type.HTTP)
                        .scheme("bearer")
                        .bearerFormat("JWT")))
            .addSecurityItem(new SecurityRequirement()
                .addList("bearer-jwt"));
    }
}
\end{lstlisting}

\textbf{Acceso a Swagger UI:} tras reiniciar la app, abrir
\url{http://localhost:8080/swagger-ui.html}. Aparece la documentación
interactiva con todos los endpoints, sus parámetros, y un botón
\emph{Try it out} para probarlos en vivo.
\end{cajaejemplo}

\section{Autenticación de APIs con JWT}

JWT (JSON Web Token, RFC 7519) es el estándar dominante de
autenticación en APIs REST modernas. Permite que el servidor
verifique la identidad del cliente sin guardar estado: la firma
criptográfica del token es prueba suficiente.


\subsection{Anatomía de un JWT}

Un JSON Web Token (RFC 7519 \cite{jones2015jwt}) consta de tres
partes separadas por puntos (Listado~\ref{lst:15.10}):

\begin{lstlisting}[style=shell,numbers=none,caption={Ejemplo de Shell},label=lst:15.10]
header.payload.signature
\end{lstlisting}

\textbf{Header (Base64URL) (Listado~\ref{lst:15.11}):}
\begin{lstlisting}[style=js,numbers=none,caption={Ejemplo de JavaScript},label=lst:15.11]
{ "alg": "HS256", "typ": "JWT" }
\end{lstlisting}

\textbf{Payload (Base64URL) (Listado~\ref{lst:15.12}):}
\begin{lstlisting}[style=js,numbers=none,caption={Ejemplo de JavaScript},label=lst:15.12]
{
  "sub":   "42",
  "iss":   "uteq.edu.ec",
  "iat":   1747000000,
  "exp":   1747003600,
  "rol":   "ROLE_ADMIN",
  "email": "admin@uteq.edu.ec",
  "jti":   "a1b2c3d4-e5f6-7890-1234-567890abcdef"
}
\end{lstlisting}

\textbf{Signature:} HMAC-SHA256 (o RSA/ECDSA) sobre
\texttt{base64(header).base64(payload)} con un secreto.

\subsection{Flags estándar del payload (\emph{claims})}

El payload de un JWT puede contener cualquier dato (claims),
pero hay un conjunto estandarizado de claims que el estudiante debe
incluir siempre. La lista siguiente los recoge.


\begin{itemize}
\item \texttt{iss} (issuer): emisor del token.
\item \texttt{sub} (subject): identificador del usuario.
\item \texttt{aud} (audience): receptor previsto.
\item \texttt{exp} (expiration): timestamp de expiración (Unix).
\item \texttt{nbf} (not before): no usable antes de.
\item \texttt{iat} (issued at): timestamp de emisión.
\item \texttt{jti} (JWT ID): identificador único para revocación.
\end{itemize}

\subsection{JWT en el PFC: configuración Spring Security}

La configuración de JWT en Spring Boot requiere tres piezas:
un filtro que valide el token, una configuración de seguridad que lo
declare, y una clase de utilidad que firme y verifique. El siguiente
ejemplo las muestra.


\begin{cajaejemplo}{Ejemplo 15.6 --- Generación y validación JWT en Spring Boot 4}

\textbf{Enunciado.} Implementar generación y validación de JWT en Spring Boot 4 con \texttt{jjwt}: clase utilitaria, filtro de seguridad y configuración Spring Security.

\textbf{Generador (servicio) (Listado~\ref{lst:15.13}):}
\begin{lstlisting}[style=java,numbers=none,caption={Ejemplo de Java},label=lst:15.13]
import io.jsonwebtoken.Jwts; // importacion de clase externa
import io.jsonwebtoken.security.Keys; // importacion de clase externa
import javax.crypto.SecretKey; // importacion de clase externa
import java.util.Date; // importacion de clase externa
import java.util.UUID; // importacion de clase externa

@Service // capa de servicio (logica de negocio)
public class JwtService { // clase publica
    private final SecretKey key; // inmutable
    private static final long DURACION_MS = 3600 * 1000; // 1h

    // metodo publico
    // metodo publico
    public JwtService(@Value("${jwt.secret}") String secreto) {
        this.key = Keys.hmacShaKeyFor(secreto.getBytes());
    }

    // metodo publico
    // metodo publico
    public String generar(String userId, String email, String rol) {
        Date ahora = new Date();
        return Jwts.builder() // retorna
            .subject(userId)
            .issuer("uteq.edu.ec")
            .claim("email", email)
            .claim("rol",   rol)
            .id(UUID.randomUUID().toString())
            .issuedAt(ahora)
            .expiration(new Date(ahora.getTime() + DURACION_MS))
            .signWith(key)
            .compact();
    }

    public Claims parsear(String token) { // metodo publico
        return Jwts.parser() // retorna
            .verifyWith(key)
            .build()
            .parseSignedClaims(token)
            .getPayload();
    }
}
\end{lstlisting}

\textbf{Cookie HttpOnly tras login (Listado~\ref{lst:15.14}):}
\begin{lstlisting}[style=java,numbers=none,caption={Ejemplo de Java},label=lst:15.14]
@PostMapping("/api/auth/login") // endpoint HTTP POST (crear)
public ResponseEntity<?> login(@RequestBody LoginDto req, // metodo publico
                               HttpServletResponse res) {
    Usuario u = autenticar(req.email(), req.password());
    String token = jwtService.generar(
        u.getId().toString(), u.getEmail(), u.getRol());

    Cookie cookie = new Cookie("ACCESS_TOKEN", token);
    cookie.setHttpOnly(true);
    cookie.setSecure(true);
    cookie.setAttribute("SameSite", "Strict");
    cookie.setMaxAge(3600);
    cookie.setPath("/");
    res.addCookie(cookie);

    // retorna
    // retorna
    return ResponseEntity.ok(Map.of("usuario", new UsuarioDto(u)));
}
\end{lstlisting}
\end{cajaejemplo}

\section{Práctica integrada: API REST consumiendo servicio externo}

La práctica integrada construye una API REST en Spring
Boot que consume un servicio externo (OpenWeatherMap), cachea las
respuestas en Redis y maneja errores correctamente.


\begin{cajaejercicio}{Ejercicio 15.1 --- API gateway con cache-aside (Sumativa)}

\textbf{Objetivo:} construir una API REST propia que consuma una API
pública externa, aplicando cache-aside con Redis para reducir
latencia y dependencia del servicio externo.

\subsubsection*{IDE: IntelliJ IDEA Ultimate}

La práctica usa IntelliJ IDEA Ultimate. Antes de empezar,
confirma que tienes JDK 21, Maven 3.9 y Docker instalados.


\subsubsection*{Paso 1 --- Levantar Redis en Docker}

Para cachear las respuestas necesitamos Redis. Levantarlo
con Docker es lo más rápido. V\'ease el Listado~\ref{lst:15.15}.


\begin{lstlisting}[style=shell,numbers=none,caption={Ejemplo de Shell},label=lst:15.15]
docker run -d --name redis-gw -p 6379:6379 redis:7-alpine # arranca un contenedor
docker exec -it redis-gw redis-cli ping # PONG
\end{lstlisting}

\subsubsection*{Paso 2 --- Crear el proyecto Spring Boot}

\texttt{File $\rightarrow$ New $\rightarrow$ Project $\rightarrow$
Spring Boot}. Configurar:
\begin{itemize}
\item Java 21, Maven, Spring Boot 4.x.
\item Group: \texttt{ec.edu.uteq}; Artifact: \texttt{api-gateway}.
\item Dependencies: \emph{Spring Web}, \emph{Spring Data Redis},
\emph{Spring Boot DevTools}, \emph{Lombok}, \emph{Validation},
\emph{springdoc-openapi-starter-webmvc-ui} (manual en pom.xml).
\end{itemize}

\subsubsection*{Paso 3 --- Configurar \texttt{application.yml}}

El archivo \texttt{application.yml} configura la API key
del servicio externo, los parámetros de Redis y el timeout HTTP. V\'ease el Listado~\ref{lst:15.16}.


\begin{lstlisting}[style=yaml,numbers=none,caption={Ejemplo de YAML},label=lst:15.16]
spring: # configuracion de Spring
  data:
    redis: # configuracion del cliente Redis
      host: localhost # host del servicio
      port: 6379 # puerto donde escucha la aplicacion
      timeout: 2s
external:
  api:
    url: https://jsonplaceholder.typicode.com
springdoc: # configuracion de OpenAPI/Swagger
  swagger-ui: # configuracion de Swagger UI
    path: /api/docs # ruta donde se expone
\end{lstlisting}

\subsubsection*{Paso 4 --- Cliente HTTP con WebClient}

Spring Boot 4 introduce \texttt{WebClient} como reemplazo
moderno del antiguo \texttt{RestTemplate}: API reactiva, no
bloqueante, integrada con \texttt{Mono}/\texttt{Flux}. V\'ease el Listado~\ref{lst:15.17}.


\begin{lstlisting}[style=java,caption={ExternalApiClient.java},label=lst:15.17]
package ec.edu.uteq.gateway.client; // paquete del archivo

import com.fasterxml.jackson.databind.JsonNode; // importacion de clase externa
import org.springframework.beans.factory.annotation.Value; // importacion de clase externa
import org.springframework.stereotype.Component; // importacion de clase externa
import org.springframework.web.client.RestClient; // importacion de clase externa
import java.time.Duration; // importacion de clase externa
import java.util.List; // importacion de clase externa

@Component // bean Spring generico
public class ExternalApiClient { // clase publica

    private final RestClient http; // inmutable

    // metodo publico
    // metodo publico
    public ExternalApiClient(@Value("${external.api.url}") String base) {
        this.http = RestClient.builder()
            .baseUrl(base)
            .build();
    }

    public List<JsonNode> listarPosts() { // metodo publico
        return http.get().uri("/posts").retrieve() // retorna
            .body(new org.springframework.core.ParameterizedTypeReference<>() {});
    }

    public JsonNode obtenerPost(int id) { // metodo publico
        // retorna
        // retorna
        return http.get().uri("/posts/{id}", id).retrieve().body(JsonNode.class);
    }
}
\end{lstlisting}

\subsubsection*{Paso 5 --- Servicio con cache-aside}

El servicio que consume la API externa debe aplicar
cache-aside: consultar Redis primero, llamar al servicio solo si
falla la caché, y persistir el resultado con TTL. V\'ease el Listado~\ref{lst:15.18}.


\begin{lstlisting}[style=java,caption={PostService.java},label=lst:15.18]
package ec.edu.uteq.gateway.service; // paquete del archivo

import com.fasterxml.jackson.databind.ObjectMapper; // importacion de clase externa
import com.fasterxml.jackson.databind.JsonNode; // importacion de clase externa
import ec.edu.uteq.gateway.client.ExternalApiClient; // importacion de clase externa
import lombok.RequiredArgsConstructor; // importacion de clase externa
import lombok.extern.slf4j.Slf4j; // importacion de clase externa
// importacion de clase externa
// importacion de clase externa
import org.springframework.data.redis.core.StringRedisTemplate;
import org.springframework.stereotype.Service; // importacion de clase externa
import java.time.Duration; // importacion de clase externa
import java.util.*; // importacion de clase externa

@Service // capa de servicio (logica de negocio)
@Slf4j
@RequiredArgsConstructor
public class PostService { // clase publica

    private final ExternalApiClient client; // inmutable
    private final StringRedisTemplate redis; // inmutable
    private final ObjectMapper json; // inmutable

    // campo privado
    // campo privado
    private static final Duration TTL = Duration.ofSeconds(60);

    public Map<String,Object> listar() throws Exception { // metodo publico
        String cached = redis.opsForValue().get("posts:all");
        if (cached != null) { // condicional
            log.debug("Cache HIT");
            return Map.of("source","cache", // retorna
                          "data", json.readTree(cached));
        }
        log.debug("Cache MISS - consultando externa");
        List<JsonNode> data = client.listarPosts();
        redis.opsForValue().set("posts:all",
            json.writeValueAsString(data), TTL);
        return Map.of("source", "api", "data", data); // retorna
    }

    // metodo publico
    // metodo publico
    public Map<String,Object> obtener(int id) throws Exception {
        String key = "posts:" + id;
        String cached = redis.opsForValue().get(key);
        if (cached != null) { // condicional
            return Map.of("source","cache", // retorna
                          "data", json.readTree(cached));
        }
        JsonNode data = client.obtenerPost(id);
        redis.opsForValue().set(key, data.toString(), TTL);
        return Map.of("source","api", "data", data); // retorna
    }
}
\end{lstlisting}

\subsubsection*{Paso 6 --- Controller con manejo robusto de errores}

El controlador expone los endpoints REST y maneja los
errores devolviendo códigos HTTP semánticamente correctos con
ProblemDetails (RFC 7807). V\'ease el Listado~\ref{lst:15.19}.


\begin{lstlisting}[style=java,caption={PostController.java},label=lst:15.19]
package ec.edu.uteq.gateway.controller; // paquete del archivo

import ec.edu.uteq.gateway.service.PostService; // importacion de clase externa
import io.swagger.v3.oas.annotations.Operation; // importacion de clase externa
import io.swagger.v3.oas.annotations.tags.Tag; // importacion de clase externa
import lombok.RequiredArgsConstructor; // importacion de clase externa
import org.springframework.http.HttpStatus; // importacion de clase externa
import org.springframework.http.ProblemDetail; // importacion de clase externa
import org.springframework.http.ResponseEntity; // importacion de clase externa
import org.springframework.web.bind.annotation.*; // importacion de clase externa
import org.springframework.web.client.RestClientException; // importacion de clase externa
import java.net.URI; // importacion de clase externa
import java.util.Map; // importacion de clase externa

@RestController // expone endpoints REST con JSON
@RequestMapping("/api/v1/publicaciones") // ruta base del controlador
@RequiredArgsConstructor
@Tag(name = "Publicaciones", // agrupa endpoints en Swagger UI
     description = "Gateway que cachea respuestas externas")
public class PostController { // clase publica

    private final PostService service; // inmutable

    @GetMapping // endpoint HTTP GET
    // OpenAPI
    // OpenAPI
    @Operation(summary = "Lista de publicaciones (con cache 60s)")
    public ResponseEntity<?> listar() { // metodo publico
        try {
            return ResponseEntity.ok(service.listar()); // retorna
        } catch (RestClientException ex) {
            // retorna
            // retorna
            return error(503, "Servicio externo no disponible",
                         ex.getMessage());
        } catch (Exception ex) {
            return error(502, "Respuesta externa invalida", // retorna
                         ex.getMessage());
        }
    }

    @GetMapping("/{id}") // endpoint HTTP GET
    @Operation(summary = "Una publicacion por ID") // OpenAPI
    // metodo publico
    // metodo publico
    public ResponseEntity<?> obtener(@PathVariable int id) {
        try {
            return ResponseEntity.ok(service.obtener(id)); // retorna
        } catch (RestClientException ex) {
            // condicional
            // condicional
            if (ex.getMessage() != null && ex.getMessage().contains("404")) {
                // retorna
                // retorna
                return error(404, "Publicacion no encontrada",
                             "ID: " + id);
            }
            // retorna
            // retorna
            return error(503, "Servicio externo no disponible",
                         ex.getMessage());
        } catch (Exception ex) {
            return error(502, "Respuesta externa invalida", // retorna
                         ex.getMessage());
        }
    }

    // campo privado
    // campo privado
    private ResponseEntity<ProblemDetail> error(int status, String titulo,
                                                String detalle) {
        ProblemDetail pd = ProblemDetail.forStatusAndDetail(
            HttpStatus.valueOf(status), detalle);
        pd.setTitle(titulo);
        pd.setType(URI.create("https://uteq.edu.ec/errors/gateway"));
        return ResponseEntity.status(status).body(pd); // retorna
    }
}
\end{lstlisting}

\subsubsection*{Paso 7 --- Probar con Postman}

Con la API corriendo, el siguiente paso es probarla con
Postman cubriendo casos de éxito y de error para validar la
robustez.


\begin{enumerate}
\item Arrancar la app: \texttt{Shift+F10}.
\item En IntelliJ Ultimate, abrir el panel \emph{HTTP Client}
(\texttt{Tools $\rightarrow$ HTTP Client $\rightarrow$ Create Request}),
o usar Postman externo.

\item Crear archivo \texttt{requests.http}:
\begin{lstlisting}[style=shell,numbers=none,caption={Ejemplo de Shell},label=lst:15.20]
### Listar publicaciones (primera vez: cache MISS)
GET http://localhost:8080/api/v1/publicaciones
Accept: application/json

### Listar otra vez (cache HIT)
GET http://localhost:8080/api/v1/publicaciones
Accept: application/json

### Obtener una
GET http://localhost:8080/api/v1/publicaciones/1
Accept: application/json

### Probar caso 404
GET http://localhost:8080/api/v1/publicaciones/9999
Accept: application/json
\end{lstlisting}

\item Ejecutar cada petición pulsando el botón verde junto a cada
\texttt{\#\#\#}.

\item Verificar:
\begin{itemize}
\item Primera llamada: campo \texttt{source: \char34{}api\char34{}}.
\item Llamada inmediata posterior: \texttt{source: \char34{}cache\char34{}}.
\item Llamada con ID inexistente: 404 con ProblemDetails.
\end{itemize}
\end{enumerate}

\subsubsection*{Paso 8 --- Verificar Swagger UI}

Abrir \url{http://localhost:8080/api/docs}. Aparece la documentación
interactiva. Probar el endpoint \texttt{GET /api/v1/publicaciones/\{id\}}
desde la interfaz Swagger.

\subsubsection*{Paso 9 --- Inspeccionar el cache en Redis}

Después de las pruebas, inspeccionar el contenido de Redis
con \texttt{redis-cli} permite verificar que las claves se están
generando correctamente con sus TTL. V\'ease el Listado~\ref{lst:15.21}.


\begin{lstlisting}[style=shell,numbers=none,caption={Ejemplo de Shell},label=lst:15.21]
docker exec -it redis-gw redis-cli # ejecuta comando en contenedor activo
> KEYS posts:*
1) "posts:all"
2) "posts:1"
> TTL posts:all
(integer) 47 # segundos restantes hasta expirar
> GET posts:1
"{\"userId\":1,\"id\":1,\"title\":\"...\",\"body\":\"...\"}"
> exit
\end{lstlisting}

\subsubsection*{Resultado esperado}

API gateway funcional que:
\begin{itemize}
\item Cachea la respuesta externa por 60 segundos.
\item Identifica explícitamente la fuente (\texttt{cache} o
\texttt{api}).
\item Maneja 4 escenarios de error con ProblemDetails RFC 7807.
\item Está documentada con Swagger UI en \texttt{/api/docs}.
\end{itemize}

\textbf{Esta es la \emph{Sumativa GA-Práctica de la semana 15}.}
\end{cajaejercicio}

\section{Práctica integrada: PFC Entrega 2 (Semana 15)}

La Entrega 2 del PFC consolida los avances de las semanas
10--15: gestión de estado, acceso a datos con ORM, arquitectura
documentada y API REST autenticada.


\begin{cajaejercicio}{Ejercicio 15.2 --- PFC Entrega 2 final --- API REST + Angular operacional}

Esta es la última entrega parcial del PFC antes de la Entrega Final
de la semana 16. Cierra todas las funcionalidades del backend con su
API REST documentada y el frontend Angular completamente operacional.

\subsubsection*{Requisitos cubiertos en esta entrega}

La Entrega 2 cubre un alcance específico que se evalúa
contra rúbrica. La lista siguiente lo detalla.


\begin{itemize}
\item API REST con todos los endpoints CRUD de las entidades del
dominio.
\item Autenticación JWT con cookie HttpOnly funcionando.
\item OpenAPI 3.0 con Swagger UI accesible en \texttt{/api/docs}.
\item Frontend Angular con login, logout y CRUD de al menos una
entidad principal.
\item Manejo de errores estandarizado con ProblemDetails (RFC 7807).
\item Cache Redis en al menos un endpoint de listado.
\item Logs estructurados en backend (SLF4J + Logback).
\item Mínimo 5 pruebas unitarias con JUnit 5 y MockMvc.
\end{itemize}

\subsubsection*{Entregables}

El equipo entrega un conjunto de artefactos consolidados a
través del repositorio Git y formularios institucionales.


\begin{enumerate}
\item Repositorio Git con tag \texttt{v0.7.0}.
\item PDF de informe técnico (15-25 páginas).
\item Demostración funcional con \texttt{docker compose up}.
\item Colección Postman exportada con todos los endpoints.
\item Captura del reporte de pruebas (JaCoCo).
\end{enumerate}

\textbf{Esta es la \emph{Sumativa TA --- PFC Entrega 2 final}}.
\end{cajaejercicio}

\section{Buenas prácticas profesionales en APIs REST}

La caja siguiente compila las reglas profesionales para
escribir APIs REST mantenibles, evolucionables y bien documentadas.


\begin{cajabuenas}{Veinte reglas profesionales para APIs REST 2026}
\begin{enumerate}
\item Versionar la API: \texttt{/api/v1/...} en URL o cabecera
\texttt{Accept: application/vnd.uteq.v1+json}.
\item Sustantivos en plural; sin verbos en URL.
\item Códigos HTTP semánticamente correctos.
\item ProblemDetails RFC 7807 para errores.
\item Validación con Bean Validation; nunca confiar en cliente.
\item Paginación obligatoria en listados (max 100 items por page).
\item Filtros, ordenamiento y búsqueda por query string.
\item Idempotencia respetada en GET, PUT, DELETE.
\item HTTPS obligatorio incluso en desarrollo local con mkcert.
\item Autenticación JWT en cookie HttpOnly+Secure+SameSite=Strict.
\item Rate limiting por IP/usuario (\texttt{X-RateLimit-*}).
\item CORS explícito y restrictivo (no \texttt{*} en producción).
\item OpenAPI documentando todos los endpoints.
\item DTOs separados de entidades JPA/EF.
\item Logging estructurado en JSON para análisis.
\item Health checks (\texttt{/actuator/health} en Spring,
\texttt{/health} en .NET).
\item Cache headers correctos (\texttt{Cache-Control}, \texttt{ETag}).
\item Compresión Gzip/Brotli activa.
\item Pruebas de contrato con OpenAPI schema validation.
\item Colección Postman versionada en el repositorio.
\end{enumerate}
\end{cajabuenas}

\section{Contraste bibliográfico de los conceptos clave}

REST como estilo arquitectónico fue formalizado por
\cite{fielding2000rest} en su tesis doctoral. La literatura más
reciente ha matizado y enriquecido ese marco. \cite{richardson2007restful}
introdujo el \emph{Richardson Maturity Model} (niveles 0--3) como
guía para evaluar el grado de adopción de los principios REST en
una API real. \cite{richardson2007restful} discute con datos empíricos
qué prácticas REST se cumplen en APIs públicas: la mayoría no
alcanza el nivel 3 (HATEOAS) y muchas violan principios básicos.
\cite{richardson2007restful} provee una taxonomía de errores comunes en
diseño REST en APIs latinoamericanas; \cite{richardson2007restful}
documenta cómo OpenAPI 3.0 + Swagger UI se han convertido en el
estándar de facto para documentación viva en empresas. Sobre SOAP,
\cite{alonso2010webservice} sigue siendo la referencia técnica, y
\cite{daigneau2012service} discute patrones de diseño para
servicios que aplican tanto a SOAP como a REST.

\section{Ejercicio resuelto adicional con resultado visual}

Como cierre de la semana 15, el siguiente ejercicio
resuelto adicional construye una API REST + JWT + Swagger paso a
paso, y se complementa con un propuesto de migración SOAP $\to$ REST.


\begin{cajaejemplo}{Ejercicio resuelto 15.1 --- API REST de productos con JWT y Swagger}
\textbf{Enunciado.} Implementar una API REST de productos con
Spring Boot 4 que requiera autenticación JWT, documente sus
endpoints con OpenAPI 3.0 y exponga Swagger UI interactivo.

\textbf{Dependencias clave (\texttt{pom.xml}):}
\begin{lstlisting}[language=yaml,caption={Dependencias Maven},label=lst:15.22]
spring-boot-starter-web
spring-boot-starter-data-jpa
spring-boot-starter-security
io.jsonwebtoken:jjwt-api:0.12.5
springdoc-openapi-starter-webmvc-ui:2.5.0
\end{lstlisting}

\textbf{Configuración de SpringDoc (\texttt{application.yml}):}
\begin{lstlisting}[language=yaml,caption={Configuración de Swagger},label=lst:15.23]
springdoc: # configuracion de OpenAPI/Swagger
  api-docs: # configuracion del JSON OpenAPI
    path: /v3/api-docs # ruta donde se expone
  swagger-ui: # configuracion de Swagger UI
    path: /swagger-ui.html # ruta donde se expone
    operations-sorter: method
\end{lstlisting}

\textbf{Endpoint anotado con OpenAPI (Listado~\ref{lst:15.24}):}
\begin{lstlisting}[language=Java,caption={Endpoint documentado},label=lst:15.24]
@RestController // expone endpoints REST con JSON
@RequestMapping("/api/v1/productos") // ruta base del controlador
// agrupa endpoints en Swagger UI
// agrupa endpoints en Swagger UI
@Tag(name = "Productos", description = "Catalogo de productos")
@SecurityRequirement(name = "bearer-jwt") // declara requisito de seguridad
public class ProductoController { // clase publica

  @Operation(summary = "Lista productos paginados", // OpenAPI
    description = "Retorna lista paginada con filtro opcional por categoria")
  @ApiResponses({ // documenta posible respuesta HTTP
    // documenta posible respuesta HTTP
    // documenta posible respuesta HTTP
    @ApiResponse(responseCode = "200", description = "Exito"),
    // documenta posible respuesta HTTP
    // documenta posible respuesta HTTP
    @ApiResponse(responseCode = "401", description = "No autenticado")
  })
  @GetMapping // endpoint HTTP GET
  public Page<Producto> listar( // metodo publico
    @RequestParam(required = false) String categoria, // query param
    Pageable pageable) {
    return categoria != null // retorna
      ? service.porCategoria(categoria, pageable)
      : service.todos(pageable);
  }
}
\end{lstlisting}

\textbf{Resultado visual} de Swagger UI en
\url{http://localhost:8080/swagger-ui.html}:

\begin{center}
\fbox{\begin{minipage}{0.95\textwidth}
\vspace{0.2cm}\footnotesize
\hspace{0.2cm}\textbf{API de Productos UTEQ} \hspace{1cm} {\color{uteqverde}OAS 3.0}\\
\hspace{0.2cm}\rule{0.95\textwidth}{0.5pt}\\[0.2cm]
\hspace{0.2cm}\textbf{Servers:} \texttt{http://localhost:8080}\\
\hspace{0.2cm}\textbf{[Autorizar]} {\small (Click para configurar JWT bearer)}\\[0.3cm]
\hspace{0.2cm}\colorbox{uteqverde!20}{\strut\ \textbf{GET}\ }\ \texttt{/api/v1/productos}\ \ \ Lista productos paginados\ \ \ $\downarrow$\\
\hspace{0.2cm}\colorbox{blue!20}{\strut\ \textbf{POST}\ }\ \texttt{/api/v1/productos}\ \ \ Crea un producto\ \ \ $\downarrow$\\
\hspace{0.2cm}\colorbox{uteqverde!20}{\strut\ \textbf{GET}\ }\ \texttt{/api/v1/productos/\{id\}}\ \ \ Obtiene un producto\ \ \ $\downarrow$\\
\hspace{0.2cm}\colorbox{uteqdorado!40}{\strut\ \textbf{PUT}\ }\ \texttt{/api/v1/productos/\{id\}}\ \ \ Actualiza un producto\ \ \ $\downarrow$\\
\hspace{0.2cm}\colorbox{red!20}{\strut\ \textbf{DELETE}\ }\ \texttt{/api/v1/productos/\{id\}}\ \ \ Elimina un producto\ \ \ $\downarrow$\\
\vspace{0.1cm}
\end{minipage}}
\end{center}

\textbf{Probando un endpoint desde Swagger UI:} al hacer clic en
\emph{Try it out} sobre \texttt{GET /api/v1/productos}, ejecuta la
petición con el JWT configurado y muestra la respuesta:

\begin{center}
\fbox{\begin{minipage}{0.95\textwidth}
\vspace{0.15cm}\footnotesize\ttfamily
\hspace{0.2cm}\textbf{Response 200 OK}\\
\hspace{0.2cm}content-type: application/json\\
\hspace{0.2cm}\\
\hspace{0.2cm}\{\\
\hspace{0.4cm}"content": [\\
\hspace{0.6cm}\{ "id": 1, "nombre": "Mouse Logitech", "precio": 25.50, "stock": 100 \},\\
\hspace{0.6cm}\{ "id": 2, "nombre": "Teclado mecánico", "precio": 89.00, "stock": 45 \}\\
\hspace{0.4cm}],\\
\hspace{0.4cm}"totalElements": 2, "totalPages": 1, "number": 0\\
\hspace{0.2cm}\}\\
\vspace{0.1cm}
\end{minipage}}
\end{center}

\textbf{Discusión.} La combinación springdoc + OpenAPI 3.0 produce
documentación viva automática a partir de las anotaciones del código.
El equipo nunca tiene que mantener un archivo Swagger a mano: cada
cambio en el controlador se refleja inmediatamente en
\texttt{/swagger-ui.html}. La integración con Spring Security permite
probar los endpoints autenticados directamente desde el navegador
\cite{richardson2007restful}.
\end{cajaejemplo}

\begin{cajaejercicio}{Ejercicio propuesto 15.2 --- Migración progresiva SOAP a REST con fachada}
\textbf{Enunciado.} Un sistema legado expone un servicio SOAP de
inventario con los métodos \texttt{consultarStock(codigo)},
\texttt{reservarItem(codigo, cantidad)} y
\texttt{cancelarReserva(idReserva)}. Construir una fachada REST
moderna que exponga estos métodos como recursos y permita una
migración progresiva.

\textbf{Requisitos:}
\begin{enumerate}
\item Backend Spring Boot que use \texttt{WebServiceTemplate} para
consumir el SOAP legado.
\item Exponer endpoints REST:
\begin{itemize}
\item \texttt{GET /api/v1/inventario/\{codigo\}}.
\item \texttt{POST /api/v1/inventario/\{codigo\}/reservas} (body
con \texttt{cantidad}).
\item \texttt{DELETE /api/v1/inventario/reservas/\{idReserva\}}.
\end{itemize}
\item Mapear errores SOAP (\emph{SOAPFault}) a respuestas
ProblemDetails (RFC 7807).
\item Caché Redis sobre \texttt{consultarStock} con TTL de 30 s para
reducir carga sobre el SOAP origen.
\item Documentar con OpenAPI y exponer Swagger UI.
\item Producir una ADR-004 que justifique la decisión de mantener
SOAP por debajo en lugar de reescribir.
\end{enumerate}

\textbf{Criterios de evaluación:}
\begin{itemize}
\item Latencia p95 $<$ 100 ms en cache caliente.
\item Manejo correcto de errores: SOAP \emph{timeout} se traduce a
HTTP 504; \emph{SOAPFault de stock insuficiente} a HTTP 409 Conflict.
\item ADR completa y razonada.
\item Pruebas Postman con 8 escenarios de éxito y error.
\end{itemize}

\textbf{Lecturas recomendadas:} \cite{alonso2010webservice} para
SOAP; \cite{newman2021building} cap. 3 para patrones de migración.
\end{cajaejercicio}

% =====================================================================
%       SEMANA 16: PFC ENTREGA FINAL - PRUEBAS, COBERTURA Y SEGURIDAD
% =====================================================================
\chapter{Entrega final del PFC: pruebas, cobertura, seguridad y publicación}
\label{cap:semana16}

\epigraph{«Un programa sin pruebas no está terminado: simplemente
todavía no sabemos en qué falla.»}{--- \textit{Adaptado de Edsger W.~Dijkstra}}

\section*{Resultado de aprendizaje de la semana}
Al concluir la semana 16, el estudiante entrega el sistema PFC en
versión 1.0.0 estable, con suite de pruebas unitarias e integración
de cobertura $\geq 70\%$, validación OWASP de los seis controles
mínimos, mediciones empíricas con k6 y Lighthouse, y la documentación
técnica final completa con mínimo 15 referencias APA.

\section{Pruebas en aplicaciones web modernas}

Una aplicación sin pruebas es como un puente sin
inspecciones: puede funcionar, pero nadie sabe cuánto resistirá al
primer terremoto. Las pruebas son la red de seguridad que permite
refactorizar con confianza.


\begin{cajahistoria}{De TDD a la pirámide de pruebas --- 25 años de cultura de testing}
La cultura de testing en software fue catapultada por Kent Beck con
\emph{Test-Driven Development} (TDD) en \emph{Extreme Programming
Explained} (1999). La idea: escribir la prueba ANTES del código,
hacer que pase con la implementación más simple posible, refactorizar.
TDD se convirtió en práctica estándar en equipos ágiles.

Mike Cohn formalizó en 2009 la \textbf{pirámide de pruebas}: muchas
pruebas unitarias rápidas en la base, menos pruebas de integración en
el medio, y unas pocas pruebas end-to-end en la cima. Esta jerarquía
balancea velocidad de feedback (las unitarias corren en milisegundos)
y confianza (las E2E validan el sistema completo).

En 2026 se ha matizado el modelo. Google propone el \textbf{«modelo
honeycomb»}: priorizar pruebas de integración medias (\emph{wide unit
tests}) sobre las unitarias puras, especialmente en arquitecturas
microservicios donde un servicio aislado dice poco. Spotify
popularizó el \textbf{«trofeo de pruebas»}: análisis estático abajo,
pocas unitarias, muchas integración, pocas E2E.
\end{cajahistoria}

\subsection{Niveles de pruebas en aplicaciones web}

Existen distintos niveles de pruebas con distintos costos y
distintos tipos de defectos que detectan. La pirámide de pruebas
(Mike Cohn) los organiza visualmente.


\begin{table}[htbp]
\centering
\caption{Niveles de pruebas y tecnologías 2026.}
\footnotesize
\begin{tabularx}{\textwidth}{lXX}
\toprule
\textbf{Nivel} & \textbf{Qué prueba} & \textbf{Herramientas (Java/Spring)}\\
\midrule
Unitarias    & Una clase aislada & JUnit 5 + Mockito + AssertJ\\
Integración  & Varias clases o capas & SpringBootTest, MockMvc, Testcontainers\\
Contrato     & Compatibilidad cliente-servidor & Pact, Spring Cloud Contract\\
E2E          & Sistema completo & Cypress, Playwright, Selenium\\
Carga / rendimiento & Sistema bajo presión & k6, JMeter, Gatling\\
Seguridad    & Vulnerabilidades & OWASP ZAP, Burp Suite\\
Mutación     & Calidad de las pruebas & PIT (Pitest), Stryker\\
\bottomrule
\end{tabularx}
\end{table}

\subsection{La pirámide de pruebas en detalle}

Cada nivel de la pirámide tiene un propósito específico:
las unitarias garantizan correctitud lógica, las de integración
verifican que las piezas encajen, y las E2E confirman que la
aplicación entrega valor al usuario.


\begin{figure}[htbp]
\centering
\begin{tikzpicture}[every node/.style={font=\scriptsize,align=center}]
  % E2E - punta
  \fill[red!30] (-1.2,3) -- (1.2,3) -- (0.6,4) -- (-0.6,4) -- cycle;
  \draw[thick] (-1.2,3) -- (1.2,3) -- (0.6,4) -- (-0.6,4) -- cycle;
  \node at (0,3.5) {\textbf{E2E}\\(pocas)};

  % Integración - medio
  \fill[uteqdorado!40] (-2.4,1.5) -- (2.4,1.5) -- (1.2,3) -- (-1.2,3) -- cycle;
  \draw[thick] (-2.4,1.5) -- (2.4,1.5) -- (1.2,3) -- (-1.2,3) -- cycle;
  \node at (0,2.25) {\textbf{Integración}\\(algunas)};

  % Unitarias - base
  \fill[uteqverde!30] (-3.6,0) -- (3.6,0) -- (2.4,1.5) -- (-2.4,1.5) -- cycle;
  \draw[thick] (-3.6,0) -- (3.6,0) -- (2.4,1.5) -- (-2.4,1.5) -- cycle;
  \node at (0,0.75) {\textbf{Pruebas unitarias}\\(muchas, rápidas)};

  % Anotaciones
  \node[anchor=west,font=\tiny] at (4,3.5)
       {Costoso, lento\\Pocas, alta confianza};
  \node[anchor=west,font=\tiny] at (4,2.25)
       {Costo medio\\Confirma flujos};
  \node[anchor=west,font=\tiny] at (4,0.75)
       {Barato, rápido\\Cubre lógica};
\end{tikzpicture}
\caption{Pirámide de pruebas de Cohn: la base son muchas pruebas
unitarias rápidas; la cumbre son pocas pruebas E2E lentas pero de
alta confianza.}
\end{figure}

\section{Pruebas unitarias con JUnit 5 y Mockito}

JUnit 5 es el framework de pruebas dominante en Java desde
2017. Mockito complementa con \emph{mocking} para aislar la unidad
bajo prueba de sus colaboradores.


\begin{cajaejemplo}{Ejemplo 16.1 --- Test unitario de un servicio con mock del repositorio}

\textbf{Enunciado.} Construir una suite completa de tests para un servicio de gestión de usuarios: tests unitarios con Mockito, tests de integración con MockMvc y Testcontainers, y reporte de cobertura JaCoCo.

\begin{lstlisting}[style=java,caption={ProductoServiceTest.java},label=lst:16.1]
package ec.edu.uteq.appweb.service; // paquete del archivo

import ec.edu.uteq.appweb.model.Producto; // importacion de clase externa
import ec.edu.uteq.appweb.repository.ProductoRepository; // importacion de clase externa
import org.junit.jupiter.api.*; // importacion de clase externa
import org.mockito.*; // importacion de clase externa
import java.math.BigDecimal; // importacion de clase externa
import java.util.*; // importacion de clase externa
import static org.assertj.core.api.Assertions.*; // importacion de clase externa
import static org.mockito.Mockito.*; // importacion de clase externa
import static org.junit.jupiter.api.Assertions.*; // importacion de clase externa

class ProductoServiceTest {

    @Mock private ProductoRepository repo; // crea un mock con Mockito
    @InjectMocks private ProductoService service; // inyecta los mocks en la clase bajo prueba
    private AutoCloseable mocks; // campo privado

    @BeforeEach // se ejecuta antes de cada test
    void setUp() { mocks = MockitoAnnotations.openMocks(this); }

    @AfterEach // se ejecuta despues de cada test
    void tearDown() throws Exception { mocks.close(); }

    @Test // metodo de test JUnit
    @DisplayName("Listar productos devuelve la lista del repositorio")
    void listar_devuelveListaDelRepo() {
        // Arrange
        Producto p1 = new Producto(1L, "A", new BigDecimal("1.00"), 5);
        Producto p2 = new Producto(2L, "B", new BigDecimal("2.00"), 10);
        when(repo.findAll()).thenReturn(List.of(p1, p2)); // configura comportamiento del mock

        // Act
        var resultado = service.listar();

        // Assert
        assertThat(resultado).hasSize(2)
            .extracting(Producto::getNombre)
            .containsExactly("A", "B");
        verify(repo, times(1)).findAll(); // verifica interaccion con el mock
    }

    @Test // metodo de test JUnit
    @DisplayName("Buscar inexistente lanza NoSuchElementException")
    void buscarInexistente_lanzaExcepcion() {
        // configura comportamiento del mock
        // configura comportamiento del mock
        when(repo.findById(999L)).thenReturn(Optional.empty());
        assertThatThrownBy(() -> service.buscar(999L))
            .isInstanceOf(NoSuchElementException.class)
            .hasMessageContaining("999");
    }

    @Test // metodo de test JUnit
    @DisplayName("Guardar persiste y retorna el producto con ID")
    void guardar_persisteYRetorna() {
        Producto entrada = new Producto(null, "X",
                new BigDecimal("3.0"), 1);
        Producto guardado = new Producto(99L, "X",
                new BigDecimal("3.0"), 1);
        when(repo.save(entrada)).thenReturn(guardado); // configura comportamiento del mock

        Producto resultado = service.guardar(entrada);

        assertThat(resultado.getId()).isEqualTo(99L);
        verify(repo).save(entrada); // verifica interaccion con el mock
    }
}
\end{lstlisting}
\end{cajaejemplo}

\section{Pruebas de integración con MockMvc y Testcontainers}

Las pruebas de integración verifican que múltiples
componentes funcionen juntos. En Spring Boot, MockMvc permite
simular peticiones HTTP sin levantar un servidor real, y
Testcontainers permite usar una BD real efímera.


\begin{cajaejemplo}{Ejemplo 16.2 --- Test de integración del controller con MockMvc}

\textbf{Enunciado.} Construir un test de integración para el controller con MockMvc: simula peticiones HTTP sin levantar un servidor real, verificando códigos de estado y contenido JSON.

\begin{lstlisting}[style=java,caption={ProductoControllerIT.java},label=lst:16.2]
package ec.edu.uteq.appweb.controller; // paquete del archivo

import org.junit.jupiter.api.Test; // importacion de clase externa
// importacion de clase externa
// importacion de clase externa
import org.springframework.beans.factory.annotation.Autowired;
// importacion de clase externa
// importacion de clase externa
import org.springframework.boot.test.context.SpringBootTest;
import org.springframework.http.MediaType; // importacion de clase externa
import org.springframework.test.context.ActiveProfiles; // importacion de clase externa
import org.springframework.test.web.servlet.MockMvc; // importacion de clase externa
// importacion de clase externa
// importacion de clase externa
import org.springframework.test.web.servlet.setup.MockMvcBuilders;
// importacion de clase externa
// importacion de clase externa
import org.springframework.web.context.WebApplicationContext;
import jakarta.transaction.Transactional; // importacion de clase externa
// importacion de clase externa
// importacion de clase externa
import static org.springframework.test.web.servlet.request.MockMvcRequestBuilders.*;
// importacion de clase externa
// importacion de clase externa
import static org.springframework.test.web.servlet.result.MockMvcResultMatchers.*;

@SpringBootTest // levanta el contexto Spring para test
@ActiveProfiles("test")
@Transactional // envuelve el metodo en transaccion
class ProductoControllerIT {

    @Autowired private WebApplicationContext wac; // inyeccion de dependencias
    private MockMvc mvc; // campo privado

    @org.junit.jupiter.api.BeforeEach
    void setUp() {
        mvc = MockMvcBuilders.webAppContextSetup(wac).build();
    }

    @Test // metodo de test JUnit
    void crear_conDatosValidos_retorna201() throws Exception {
        String body = """
            {"nombre":"Test","precio":4.50,"stock":10}
            """;

        mvc.perform(post("/api/v1/productos")
                .contentType(MediaType.APPLICATION_JSON)
                .content(body))
            .andExpect(status().isCreated())
            .andExpect(header().exists("Location"))
            .andExpect(jsonPath("$.id").exists())
            .andExpect(jsonPath("$.nombre").value("Test"));
    }

    @Test // metodo de test JUnit
    void crear_conNombreInvalido_retorna422() throws Exception {
        String body = """
            {"nombre":"","precio":4.50,"stock":10}
            """;
        mvc.perform(post("/api/v1/productos")
                .contentType(MediaType.APPLICATION_JSON)
                .content(body))
            .andExpect(status().isUnprocessableEntity())
            .andExpect(jsonPath("$.errores.nombre").exists());
    }

    @Test // metodo de test JUnit
    void obtener_inexistente_retorna404() throws Exception {
        mvc.perform(get("/api/v1/productos/999999"))
            .andExpect(status().isNotFound())
            .andExpect(jsonPath("$.title").value("Recurso no encontrado"));
    }
}
\end{lstlisting}
\end{cajaejemplo}

\section{Cobertura con JaCoCo}

\textbf{JaCoCo} (Java Code Coverage) mide qué porcentaje de las
líneas y ramas del código están cubiertas por las pruebas. Es la
herramienta estándar de la industria para Java en 2026.

\begin{cajaejemplo}{Ejemplo 16.3 --- Configurar JaCoCo en Maven}

\textbf{Enunciado.} Configurar JaCoCo en un proyecto Maven para generar reportes HTML de cobertura de tests con umbrales mínimos exigibles en CI.

\textbf{En \texttt{pom.xml}:}
\begin{lstlisting}[style=xml,caption={Configuración de JaCoCo},label=lst:16.3]
<build>
  <plugins>
    <plugin>
      <groupId>org.jacoco</groupId>
      <artifactId>jacoco-maven-plugin</artifactId>
      <version>0.8.12</version>
      <executions>
        <execution>
          <goals>
            <goal>prepare-agent</goal>
          </goals>
        </execution>
        <execution>
          <id>generate-report</id>
          <phase>verify</phase>
          <goals><goal>report</goal></goals>
        </execution>
        <execution>
          <id>check-coverage</id>
          <phase>verify</phase>
          <goals><goal>check</goal></goals>
          <configuration>
            <rules>
              <rule>
                <element>BUNDLE</element>
                <limits>
                  <limit>
                    <counter>LINE</counter>
                    <value>COVEREDRATIO</value>
                    <minimum>0.70</minimum>
                  </limit>
                </limits>
              </rule>
            </rules>
          </configuration>
        </execution>
      </executions>
    </plugin>
  </plugins>
</build>
\end{lstlisting}

\textbf{Ejecutar (Listado~\ref{lst:16.4}):}
\begin{lstlisting}[style=shell,numbers=none,caption={Ejemplo de Shell},label=lst:16.4]
./mvnw clean verify
# Si la cobertura es < 70%, el build FALLA
# Reporte HTML en target/site/jacoco/index.html
\end{lstlisting}
\end{cajaejemplo}

\section{Pruebas de carga con k6}

k6 (de Grafana Labs) es una herramienta moderna de pruebas de carga
escrita en Go con scripts en JavaScript. Es la opción dominante en
2026 por su simplicidad y rendimiento.

\begin{cajaejemplo}{Ejemplo 16.4 --- Script k6 para el endpoint de productos}

\textbf{Enunciado.} Escribir un script de prueba de carga con \texttt{k6} (JavaScript) que dispare 100 usuarios virtuales contra el endpoint de productos durante 30 segundos.

\begin{lstlisting}[style=js,caption={carga-productos.js},label=lst:16.5]
import http from 'k6/http'; // importacion ES Modules
import { check, sleep } from 'k6'; // importacion ES Modules

export const options = { // exporta del modulo
  stages: [
    { duration: '30s', target: 10 }, // ramp-up a 10 VUs
    { duration: '1m',  target: 50 }, // ramp-up a 50 VUs
    { duration: '2m',  target: 50 }, // mantener 50 VUs
    { duration: '30s', target: 0 }, // ramp-down
  ],
  thresholds: {
    http_req_duration: ['p(95)<500'], // 95% < 500ms
    http_req_failed:   ['rate<0.01'], // < 1% errores
  },
};

const BASE = 'http://localhost:8080';

export default function () { // exporta del modulo
  // 80% de lecturas (listado), 20% de detalle
  if (Math.random() < 0.8) { // condicional
    // constante (no reasignable)
    // constante (no reasignable)
    const r = http.get(`${BASE}/api/v1/productos?page=1&size=20`);
    check(r, {
      '200 OK':           (r) => r.status === 200,
      'tiempo aceptable': (r) => r.timings.duration < 500,
    });
  } else {
    const id = Math.floor(Math.random() * 1000) + 1; // constante (no reasignable)
    const r = http.get(`${BASE}/api/v1/productos/${id}`); // constante (no reasignable)
    check(r, {
      'OK o 404': (r) => r.status === 200 || r.status === 404,
    });
  }
  sleep(0.5 + Math.random());
}
\end{lstlisting}

\textbf{Ejecutar (Listado~\ref{lst:16.6}):}
\begin{lstlisting}[style=shell,numbers=none,caption={Ejemplo de Shell},label=lst:16.6]
# Instalar k6 (Linux)
sudo apt install k6 # o brew install k6 en macOS

# Ejecutar la prueba
k6 run carga-productos.js

# Generar graficos con InfluxDB+Grafana
k6 run --out influxdb=http://localhost:8086/k6 carga-productos.js
\end{lstlisting}
\end{cajaejemplo}

\section{Auditoría de frontend con Lighthouse}

\textbf{Lighthouse} es la herramienta open source de Google integrada
en Chrome DevTools que audita una página web en cuatro dimensiones:
\textbf{Performance, Accessibility, Best Practices, SEO}.

\subsection{Las cuatro métricas core}

Lighthouse mide la calidad de una página web en cuatro
ejes: Performance, Accessibility, Best Practices, SEO. Las cuatro
métricas \emph{core} de Performance son las que más impacto tienen
en la experiencia real.


\begin{table}[htbp]
\centering
\caption{Métricas Core Web Vitals 2026.}
\footnotesize
\begin{tabularx}{\textwidth}{lXl}
\toprule
\textbf{Métrica} & \textbf{Significado} & \textbf{Umbral bueno}\\
\midrule
LCP (Largest Contentful Paint) & Tiempo hasta el elemento principal visible & $<$ 2.5 s\\
INP (Interaction to Next Paint) & Tiempo de respuesta a interacciones & $<$ 200 ms\\
CLS (Cumulative Layout Shift) & Inestabilidad visual durante la carga & $<$ 0.1\\
FCP (First Contentful Paint) & Primer pixel pintado & $<$ 1.8 s\\
TTFB (Time To First Byte) & Latencia del servidor & $<$ 800 ms\\
\bottomrule
\end{tabularx}
\end{table}

\subsection{Cómo ejecutar Lighthouse}

Lighthouse se ejecuta desde DevTools del navegador o desde
la línea de comandos. La forma DevTools es la más directa para una
auditoría puntual.


\begin{enumerate}
\item Abrir el sitio Angular en Chrome.
\item F12 $\rightarrow$ pestaña \emph{Lighthouse}.
\item Marcar \emph{Performance, Accessibility, Best Practices, SEO}.
\item Modo: \emph{Navigation (Default)}; Device: \emph{Mobile}.
\item Pulsar \emph{Analyze page load}.
\item En 1-2 minutos aparecen los 4 scores (0--100) con
recomendaciones concretas.
\end{enumerate}

\subsection{CLI alternativa para CI/CD}

Para integrar Lighthouse en un pipeline CI/CD, conviene
usar la CLI con configuración versionable. La herramienta \texttt{lhci}
(Lighthouse CI) está diseñada para esto. V\'ease el Listado~\ref{lst:16.7}.


\begin{lstlisting}[style=shell,numbers=none,caption={Ejemplo de Shell},label=lst:16.7]
npm install -g lighthouse # instala dependencias del package.json

# Auditar y guardar reporte HTML
lighthouse https://app.uteq.edu.ec --output=html \
  --output-path=./reporte.html

# Para CI/CD: failear si Performance < 80
lighthouse https://app.uteq.edu.ec --chrome-flags="--headless" \
  --output=json --output-path=./report.json
node -e "const r=require('./report.json');
  if(r.categories.performance.score<0.8) process.exit(1);"
\end{lstlisting}

\section{Auditoría de seguridad: 6 controles OWASP mínimos}

Antes de declarar listo el PFC, conviene ejecutar una
auditoría de seguridad mínima: seis controles OWASP que detectan
los problemas más graves antes de que lo haga un atacante.


\begin{table}[htbp]
\centering
\caption{Los seis controles mínimos OWASP a verificar en el PFC.}
\footnotesize
\begin{tabularx}{\textwidth}{lXX}
\toprule
\textbf{OWASP} & \textbf{Control} & \textbf{Cómo verificar}\\
\midrule
A01 & Acceso a recurso ajeno bloqueado & Login user A; intentar GET /api/usuarios/B/datos $\to$ 403\\
A02 & TLS 1.2/1.3, sin texto plano & \texttt{curl -v https://...} ver \texttt{TLSv1.3}\\
A03 & SQL Injection prevenida & Probar payload \texttt{' OR '1'='1} $\to$ 422\\
A05 & Cabeceras de seguridad activas & \texttt{curl -I} y verificar HSTS, CSP, X-Frame-Options\\
A07 & Rate limiting + JWT cookie HttpOnly & 6 intentos login $\to$ 429; cookie HttpOnly visible en DevTools\\
A09 & Logs de eventos de seguridad & Login fallido aparece en log con IP y timestamp\\
\bottomrule
\end{tabularx}
\end{table}

\section{Práctica integrada: PFC Entrega Final v1.0.0}

La Entrega Final del PFC consolida 14 semanas de trabajo
en una aplicación profesional. Se evalúa contra una rúbrica
detallada conocida de antemano.


\begin{cajaejercicio}{Ejercicio 16.1 --- PFC Entrega Final --- Sumativa principal del semestre}

Esta es la \textbf{Sumativa Principal Final} del PFC. Concentra el
trabajo de las 16 semanas previas en una entrega de calidad
profesional.

\subsubsection*{Composición de la nota (100 puntos)}

La nota de la Entrega Final se compone de cinco bloques
ponderados que el estudiante debe atender en conjunto. La lista
siguiente los detalla.


\begin{tabularx}{\textwidth}{lXl}
\toprule
\textbf{Componente} & \textbf{Criterio} & \textbf{Pts}\\
\midrule
Demo en vivo & \texttt{docker compose up} arranca todo; CRUD funcional & 20 \\
Completitud funcional & Todas las funcionalidades en estado «Completo» & 15 \\
Pruebas y cobertura & JaCoCo $\geq$ 70\%; k6 con 50 VUs; Lighthouse $\geq$ 80 & 15 \\
Comparativas técnicas & Tablas REST/SOAP, Angular/React, etc.\ con datos reales & 10 \\
Reflexión sobre tendencias & $\geq$ 400 palabras (Jamstack, PWA, IA) & 10 \\
Informe técnico & Abstract bilingüe; 15+ refs APA; conclusiones por objetivo & 15 \\
Seguridad OWASP & 6 controles verificados con evidencia & 10 \\
Repositorio final & Tag v1.0.0; commits de todos; CI verde & 5 \\
\bottomrule
\end{tabularx}

\subsubsection*{Lista de verificación obligatoria}

Antes de la sesión de laboratorio del lunes, cada equipo debe poder
marcar TODOS estos puntos:

\begin{enumerate}
\item \texttt{docker compose up} arranca todo en menos de 2 minutos;
todos los servicios \texttt{healthy}.
\item El docente puede hacer login, navegar el dashboard, crear,
editar y eliminar registros desde Angular.
\item Swagger UI accesible en \texttt{/api/docs} con todos los
endpoints documentados.
\item Abstract bilingüe (200--250 palabras en español e igual en
inglés).
\item Tabla de inventario completa; las parciales/no implementadas
marcadas.
\item Diagrama C4 Nivel 2 final en PNG, refleja despliegue real.
\item Diagrama de clases UML refactorizado en PNG/SVG ($\geq$ 300 dpi).
\item DER de la BD generado con pgAdmin 4 ERD Tool.
\item Tabla comparativa REST vs SOAP completa con ejemplo ecuatoriano.
\item Sección de tendencias (Jamstack, PWA, IA) $\geq$ 400 palabras
con posición crítica propia.
\item Reporte JaCoCo cobertura $\geq$ 70\%; captura HTML incluida.
\item Prueba de carga k6 con 50 VUs documentada.
\item Informe Lighthouse con los 4 scores del frontend Angular.
\item Tabla de seguridad OWASP con 6 controles verificados con
evidencia.
\item Conclusiones --- un párrafo por objetivo específico.
\item Reflexión crítica sobre el proceso (300--400 palabras).
\item Trabajo futuro: 5 mejoras con justificación y estimación.
\item Mínimo 15 referencias APA, $\geq$ 5 indexadas (JCR/SJR).
\item Historial Git con commits de todos los integrantes.
\item Tag \texttt{v1.0.0} en el repositorio.
\item Pipeline GitHub Actions en verde.
\item Colección Postman exportada en \texttt{docs/postman\_collection.json}.
\item PDF nombrado \texttt{PFC\_EntregaFinal\_NombreEquipo.pdf}.
\end{enumerate}

\subsubsection*{IDE y herramientas}

Esta entrega exige varias herramientas conectadas. La lista
siguiente confirma cuáles deben estar instaladas antes de la
auditoría final.


\begin{itemize}
\item IntelliJ IDEA Ultimate (backend Spring Boot).
\item VS Code con Angular Language Service (frontend Angular).
\item Docker Desktop para orquestación local.
\item pgAdmin 4 para administración de BD y generación del DER.
\item Postman para colecciones de prueba.
\item k6 para pruebas de carga.
\item Chrome DevTools + Lighthouse para auditoría frontend.
\end{itemize}

\subsubsection*{Sesión de laboratorio: viernes de la semana 16}

Cada equipo demuestra al docente:
\begin{enumerate}
\item Desde una clonación fresca del repositorio
(\texttt{git clone}), \texttt{docker compose up -d} y todo
funciona.
\item Login en Angular con usuario admin.
\item Crear, editar y eliminar al menos un registro de cada entidad
principal.
\item Mostrar reportes de pruebas (HTML JaCoCo + Lighthouse + k6).
\item Mostrar evidencia de los 6 controles OWASP.
\end{enumerate}

\textbf{Esta es la \emph{Sumativa Principal Final --- PFC v1.0.0}}.
\end{cajaejercicio}

\section{Buenas prácticas profesionales en pruebas y entregas}

La caja siguiente compila las reglas profesionales de
pruebas y entregas: principios que separan al equipo que aprueba
con tranquilidad del que sufre la noche anterior.


\begin{cajabuenas}{Quince reglas profesionales 2026}
\begin{enumerate}
\item TDD cuando el dominio sea claro; tests-first para los servicios
clave.
\item Pirámide de pruebas: muchas unitarias, algunas integración,
pocas E2E.
\item Cobertura $\geq$ 70\% en líneas para código de negocio.
\item Tests con nombres descriptivos: \emph{«crear con datos
inválidos retorna 422»}.
\item Patrón AAA (Arrange-Act-Assert) en cada test.
\item Independencia: ningún test depende de otro.
\item Datos de prueba con fixtures o factories, no hardcodeados.
\item Testcontainers para BD real en pruebas de integración.
\item Pruebas de seguridad automatizadas en el pipeline CI.
\item k6 o Gatling integrados al pipeline para detectar regresiones.
\item Lighthouse en CI: failear si performance cae bajo umbral.
\item Versionado semántico (SemVer) en tags Git.
\item Conventional Commits para historial limpio.
\item ChangeLog generado automáticamente desde commits.
\item Releases anotadas en GitHub con resumen ejecutivo.
\end{enumerate}
\end{cajabuenas}

\section{Contraste bibliográfico de los conceptos clave}

La literatura técnica sobre pruebas ha evolucionado desde el modelo
clásico de \cite{martin2008clean} ---que defiende TDD como
disciplina profesional--- hasta las visiones más empíricas actuales.
\cite{ortega2022securecoding} documenta cómo la integración de
pruebas de seguridad en el pipeline CI/CD reduce vulnerabilidades en
producción en un 60\%. \cite{verizon2025dbir} en su informe anual
sigue mostrando que las aplicaciones web son el vector de ataque más
común. Sobre Lighthouse y \emph{Core Web Vitals}, los reportes
oficiales de Google evidencian correlación entre LCP $<$ 2.5 s y
mejora del 24\% en tasa de conversión en e-commerce. La prueba de
carga con \texttt{k6} se ha vuelto estándar; \cite{kleppmann2017designing}
muestra que sin pruebas de carga el 40\% de los despliegues con
nuevos endpoints sufren incidentes en las primeras 48 horas.

\section{Ejercicio resuelto adicional con resultado visual}

Como cierre, el siguiente ejercicio resuelto adicional
implementa una suite completa de pruebas (unitarias + integración +
cobertura JaCoCo) y se complementa con un propuesto.


\begin{cajaejemplo}{Ejercicio resuelto 16.1 --- Suite completa de tests para servicio de usuarios}
\textbf{Enunciado.} Implementar y ejecutar pruebas unitarias y de
integración sobre un servicio de gestión de usuarios en Spring Boot,
medir cobertura con JaCoCo y producir el reporte HTML.

\textbf{Test unitario con Mockito (\texttt{UsuarioServiceTest.java}):}
\begin{lstlisting}[language=Java,caption={Test unitario con Mockito},label=lst:16.8]
@ExtendWith(MockitoExtension.class) // extension JUnit 5
class UsuarioServiceTest {

  @Mock UsuarioRepository repo; // crea un mock con Mockito
  @Mock PasswordEncoder encoder; // crea un mock con Mockito
  @InjectMocks UsuarioService service; // inyecta los mocks en la clase bajo prueba

  @Test // metodo de test JUnit
  void crear_conDatosValidos_persisteYRetorna() {
    // Arrange
    Usuario nuevo = new Usuario("ana@uteq.edu.ec", "Pass123!");
    // configura comportamiento del mock
    // configura comportamiento del mock
    when(encoder.encode("Pass123!")).thenReturn("$2a$10$hash");
    when(repo.save(any())).thenAnswer(inv -> { // configura comportamiento del mock
      Usuario u = inv.getArgument(0);
      u.setId(1L);
      return u; // retorna
    });

    // Act
    Usuario creado = service.crear(nuevo);

    // Assert
    assertEquals(1L, creado.getId()); // esperado vs obtenido
    assertEquals("$2a$10$hash", creado.getPasswordHash()); // esperado vs obtenido
    verify(repo).save(any()); // verifica interaccion con el mock
  }

  @Test // metodo de test JUnit
  void crear_conEmailDuplicado_lanzaExcepcion() {
    // configura comportamiento del mock
    // configura comportamiento del mock
    when(repo.existsByEmail("ana@uteq.edu.ec")).thenReturn(true);

    assertThrows(EmailDuplicadoException.class, // verifica que se lance una excepcion
      () -> service.crear(new Usuario("ana@uteq.edu.ec", "Pass123!")));
    verify(repo, never()).save(any()); // verifica interaccion con el mock
  }
}
\end{lstlisting}

\textbf{Test de integración con MockMvc (\texttt{UsuarioControllerIT.java}):}
\begin{lstlisting}[language=Java,caption={Integración con MockMvc},label=lst:16.9]
@SpringBootTest // levanta el contexto Spring para test
@AutoConfigureMockMvc // configura MockMvc automaticamente
@Testcontainers // metodo de test JUnit
class UsuarioControllerIT {

  @Container // contenedor Docker para tests (Testcontainers)
  static PostgreSQLContainer<?> postgres =
    new PostgreSQLContainer<>("postgres:16-alpine");

  @DynamicPropertySource
  static void props(DynamicPropertyRegistry r) {
    r.add("spring.datasource.url", postgres::getJdbcUrl);
    r.add("spring.datasource.username", postgres::getUsername);
    r.add("spring.datasource.password", postgres::getPassword);
  }

  @Autowired MockMvc mvc; // inyeccion de dependencias

  @Test // metodo de test JUnit
  void POST_usuarios_conDatosValidos_retorna201() throws Exception {
    String body = """
      {"email":"juan@uteq.edu.ec","password":"Segura123!"}""";

    mvc.perform(post("/api/v1/usuarios")
        .contentType(MediaType.APPLICATION_JSON)
        .content(body))
      .andExpect(status().isCreated())
      .andExpect(jsonPath("$.id").isNumber())
      .andExpect(jsonPath("$.email").value("juan@uteq.edu.ec"));
  }
}
\end{lstlisting}

\textbf{Ejecución} \texttt{mvn clean test jacoco:report}:

\begin{center}
\fbox{\begin{minipage}{0.95\textwidth}
\vspace{0.15cm}\footnotesize\ttfamily
\hspace{0.2cm}[INFO] Running com.uteq.UsuarioServiceTest\\
\hspace{0.2cm}[INFO] Tests run: 8, Failures: 0, Errors: 0, Skipped: 0\\[0.15cm]
\hspace{0.2cm}[INFO] Running com.uteq.UsuarioControllerIT\\
\hspace{0.2cm}[INFO] Tests run: 5, Failures: 0, Errors: 0, Skipped: 0\\[0.15cm]
\hspace{0.2cm}[INFO] Results:\\
\hspace{0.2cm}[INFO] Tests run: 13, Failures: 0, Errors: 0, Skipped: 0\\
\hspace{0.2cm}[INFO] BUILD SUCCESS\\
\hspace{0.2cm}[INFO] Total time: 18.342 s\\
\vspace{0.1cm}
\end{minipage}}
\end{center}

\textbf{Reporte HTML de JaCoCo} (\texttt{target/site/jacoco/index.html}):

\begin{center}
\fbox{\begin{minipage}{0.95\textwidth}
\vspace{0.15cm}\footnotesize\centering
\textbf{JaCoCo Coverage Report --- UsuariosApp v1.0.0}\\[0.2cm]
\begin{tabular}{|l|c|c|c|}
\hline
\textbf{Elemento} & \textbf{Líneas \%} & \textbf{Branches \%} & \textbf{Métodos \%} \\
\hline
com.uteq.usuarios.service & 92\% & 88\% & 100\% \\
com.uteq.usuarios.controller & 85\% & 80\% & 100\% \\
com.uteq.usuarios.security & 78\% & 75\% & 90\% \\
\hline
\textbf{Total} & \textbf{86\%} & \textbf{82\%} & \textbf{97\%} \\
\hline
\end{tabular}\\[0.1cm]
\end{minipage}}
\end{center}

\textbf{Discusión.} La pirámide de pruebas se materializa: 8 tests
unitarios rápidos (sin Spring) y 5 de integración con
testcontainers. La cobertura del 86\% supera el umbral $\geq 70\%$
exigido para la entrega del PFC. Los \emph{branches} (caminos
condicionales) son siempre menores que las líneas porque algunas
ramas raras (como excepciones de BD) no se cubren
\cite{martin2008clean}.
\end{cajaejemplo}

\begin{cajaejercicio}{Ejercicio propuesto 16.2 --- Pipeline CI con GitHub Actions y Lighthouse}
\textbf{Enunciado.} Configurar un pipeline GitHub Actions que se
ejecute en cada \emph{push} a la rama \texttt{main} y realice:
\begin{enumerate}
\item Build con Maven (\texttt{mvn clean verify}).
\item Tests unitarios e integración (con servicios PostgreSQL y
Redis levantados como servicios del pipeline).
\item Reporte de cobertura JaCoCo subido como artifact.
\item Build de imagen Docker.
\item Pruebas de carga k6 contra el container recién construido
(50 VUs, 30 s).
\item Auditoría Lighthouse del frontend (CLI con \texttt{npx lhci
autorun}), umbral mínimo Performance 80, Accessibility 90.
\item Si todo pasa: push de la imagen a GitHub Container Registry
con tag \texttt{v$\{$run\_number$\}$}.
\end{enumerate}

\textbf{Criterios de evaluación:}
\begin{itemize}
\item El badge de CI en el README muestra estado \emph{passing}.
\item Pull requests no se pueden mergear si el pipeline falla
(\emph{branch protection rule}).
\item El \emph{action} sube reportes (JaCoCo, k6, Lighthouse) como
artifacts navegables en la UI de GitHub.
\item Captura de pantalla del Actions con tres ejecuciones exitosas
en el reporte del PFC.
\end{itemize}

\textbf{Lecturas recomendadas:} \cite{martin2017clean} para
disciplina; \cite{ortega2022securecoding} para integración de
seguridad en CI.
\end{cajaejercicio}

% =====================================================================
%       SEMANA 17: DEFENSA DEL PFC Y EP2
% =====================================================================
\chapter{Defensa del PFC y Evaluación Parcial Segundo Corte}
\label{cap:semana17}

\epigraph{«Saber programar es importante; saber explicar lo que has
programado es lo que separa a un técnico de un ingeniero.»}{---
\textit{Adaptación de un consejo común en revisiones académicas}}

\section*{Naturaleza de esta semana}
La semana 17 cierra el segundo corte de la asignatura. Tiene dos
componentes evaluativos principales:

\begin{enumerate}
\item \textbf{Defensa oral del PFC v1.0.0}: cada equipo presenta y
defiende su sistema durante 30--40 minutos. Cada integrante debe
responder preguntas técnicas sobre cualquier parte del sistema, no
solo sobre la suya.
\item \textbf{Evaluación Parcial Segundo Corte}: examen escrito que
cubre las Unidades III y IV (semanas 10--16).
\end{enumerate}

\section{Cómo preparar la defensa oral del PFC}

La defensa oral del PFC es la culminación del curso. Cómo
se prepara la presentación es tan importante como qué contiene:
20 minutos pueden hacer brillar o hundir un proyecto que costó 14
semanas.


\subsection{Estructura recomendada de la presentación}

La presentación de defensa sigue una estructura clásica que
distribuye los 20 minutos en bloques temáticos. Respetarla evita
omitir contenido crítico y mantiene al tribunal enganchado.


\begin{cajadefinicion}{Estructura óptima 30 minutos --- 8 bloques}
\begin{enumerate}
\item \textbf{Portada y presentación del equipo} (1 min): nombres,
roles asumidos en el proyecto.
\item \textbf{Problema y contexto} (3 min): cuál es el problema real
que resuelve el sistema, a quién afecta, datos cuantitativos del
contexto ecuatoriano si aplica.
\item \textbf{Objetivos y alcance} (2 min): objetivo general, 3--5
específicos. Lo que está dentro y fuera del alcance.
\item \textbf{Arquitectura} (5 min): diagrama C4 nivel 2, ADRs
clave, justificación de las decisiones tecnológicas más importantes.
\item \textbf{Demo en vivo} (10 min): \texttt{docker compose up},
login, navegación, CRUD principal, Swagger UI, pruebas Postman,
evidencia de OWASP.
\item \textbf{Resultados de pruebas} (3 min): JaCoCo, k6, Lighthouse
con scores reales medidos.
\item \textbf{Reflexión y trabajo futuro} (3 min): qué se aprendió,
qué se haría distinto, propuestas concretas de mejora.
\item \textbf{Preguntas del tribunal} (5--10 min): cualquier integrante
puede ser interrogado sobre cualquier parte.
\end{enumerate}
\end{cajadefinicion}

\subsection{Las 20 preguntas más frecuentes del tribunal}

El tribunal hace preguntas con patrones reconocibles. La
lista siguiente reúne las 20 preguntas más frecuentes con la
respuesta esperada.


\begin{cajabuenas}{Banco de preguntas que el equipo debe poder responder}
\begin{enumerate}
\item ¿Por qué eligieron \emph{esa} base de datos y no otra?
\item ¿Cómo previene su sistema la inyección SQL?
\item ¿Qué pasa si el usuario X intenta acceder al recurso del
usuario Y?
\item ¿Qué hace su sistema cuando expira el JWT?
\item ¿Por qué cookies HttpOnly y no \texttt{localStorage}?
\item ¿Cómo escala el sistema si reciben 1000 usuarios concurrentes?
\item ¿Qué pasa si Redis se cae?
\item ¿Cómo invalidan el cache cuando cambia un dato?
\item ¿Cuál es el peor cuello de botella de su sistema?
\item ¿Cómo aseguran que los DTOs no expongan campos sensibles?
\item Muestren un test que falle si yo introduzco un bug obvio.
\item ¿Cómo hacen el deploy en producción real (no Docker Compose)?
\item ¿Cómo gestionan secretos (passwords, JWT secret)?
\item ¿Por qué eligieron Spring Boot sobre Quarkus o Micronaut?
\item ¿Cómo validan los datos de entrada en la API?
\item ¿Cómo manejan transacciones que cruzan varios servicios?
\item ¿Qué pasa si dos usuarios editan el mismo registro a la vez?
\item ¿Han verificado los 10 OWASP del Top 10?
\item Si tuvieran que reescribir el sistema desde cero, ¿qué
cambiarían?
\item ¿Cuántas líneas de código tiene el proyecto y cómo se reparten?
\end{enumerate}
\end{cajabuenas}

\subsection{Errores comunes en defensas anteriores}

Los errores que arruinan defensas no son técnicos: son de
comunicación, gestión del tiempo y actitud. La lista siguiente los
recoge para que el estudiante los evite.


\begin{cajaadvertencia}{Errores que cuestan calificación}
\begin{itemize}
\item \textbf{Decir «no sé» sin proponer cómo averiguarlo.} Nadie
sabe todo; se valora más decir «no estoy seguro, pero lo verificaría
ejecutando esto\dots».
\item \textbf{Desconocer la propia arquitectura}: «yo solo trabajé en
el frontend» NO es excusa; cualquier integrante debe poder responder
sobre cualquier parte.
\item \textbf{Demo no preparada}: aparecer al laboratorio sin haber
probado la demo en una computadora distinta de la del desarrollador.
\item \textbf{Slides con código inservible}: capturas de IDE
ilegibles. Mejor reescribir el código en la slide con sintaxis
resaltada.
\item \textbf{Tiempo mal distribuido}: 25 minutos de problema y
contexto, 5 de demo. La demo es el corazón.
\item \textbf{Equipo donde solo habla uno}: el tribunal pregunta
explícitamente al silencioso.
\item \textbf{Datos inventados en la presentación}: decir «el sistema
soporta 10 mil usuarios concurrentes» sin haberlo medido.
\end{itemize}
\end{cajaadvertencia}

\section{Estructura de la Evaluación Parcial Segundo Corte}

La EP2 cubre exclusivamente las \textbf{Unidades III y IV} (semanas
10--16). Se compone de cuatro partes equilibradas:

\begin{table}[htbp]
\centering
\caption{Composición de la Evaluación Parcial Segundo Corte.}
\footnotesize
\begin{tabularx}{\textwidth}{lXl}
\toprule
\textbf{Parte} & \textbf{Tipo} & \textbf{Pts}\\
\midrule
A & Selección múltiple y V/F sobre conceptos & 20 \\
B & Respuesta corta (3 líneas máx, definiciones) & 20 \\
C & Análisis de código (identificar errores y vulnerabilidades) & 30 \\
D & Ejercicio práctico integrador (escribir código completo) & 30 \\
\bottomrule
\end{tabularx}
\end{table}

\subsection{Temas garantizados en la EP2}

La EP2 cubre con seguridad ciertos temas que aparecen en
todas las ediciones del examen. Asegurar dominio sobre ellos
garantiza un mínimo razonable.


\begin{cajaadvertencia}{Temario completo EP2}
\textbf{Unidad III (Semanas 10--12):}
\begin{itemize}
\item Gestión de estado: cookies, sesiones, JWT --- diferencias y
cuándo usar cada una.
\item Flags de seguridad de cookies: \texttt{Secure}, \texttt{HttpOnly},
\texttt{SameSite}.
\item Objeto Request: lectura segura de parámetros, manejo de
\texttt{X-Forwarded-For}.
\item Objeto Response: códigos HTTP, cabeceras de seguridad,
ProblemDetails (RFC 7807).
\item PDO con prepared statements; prevención SQL Injection.
\item Patrón Repository.
\item Spring Data JPA (consultas derivadas, JPQL).
\item Problema N+1 y soluciones (\texttt{JOIN FETCH}, \texttt{with},
\texttt{Include}).
\item Migraciones con Flyway/Liquibase.
\item Procedimientos almacenados: cuándo usar.
\item Patrones arquitectónicos: capas, hexagonal, SOA, microservicios,
EDA, P2P.
\item ISO/IEC 25010 (atributos de calidad).
\item Modelo C4 (4 niveles).
\item ADR (Architectural Decision Records).
\end{itemize}

\textbf{Unidad IV (Semanas 13--16):}
\begin{itemize}
\item Escalabilidad vertical vs horizontal.
\item Teorema CAP de Brewer.
\item Algoritmos de balanceo de carga.
\item Patrones de caché: cache-aside, read-through, write-through,
write-behind.
\item Cache-aside con Redis y TTL.
\item Estrategias de invalidación de caché.
\item ATAM (Architecture Tradeoff Analysis Method).
\item Patrón MVC: ciclo de vida de petición Spring.
\item MVC server-side vs API REST + SPA.
\item REST: 6 restricciones de Fielding.
\item Modelo de Madurez de Richardson.
\item Verbos HTTP, idempotencia, seguridad.
\item Convenciones URI RESTful.
\item REST vs SOAP: comparativa, casos de uso.
\item OpenAPI 3.0 (Swagger).
\item Anatomía de un JWT (3 partes).
\item Pruebas: unitarias, integración, E2E (pirámide).
\item Cobertura JaCoCo $\geq$ 70\%.
\item k6 y Lighthouse.
\item OWASP Top 10:2021 (especialmente A01, A02, A03, A07).
\end{itemize}
\end{cajaadvertencia}

\subsection{Ejemplos del tipo de preguntas que aparecen}

Para preparar el examen conviene haber visto el tipo de
preguntas que se formulan. Los ejemplos siguientes son
representativos de las ediciones recientes.


\begin{cajaejemplo}{Muestra de la Parte A (selección múltiple)}
\textbf{Pregunta A1:} ¿Cuál de las siguientes secuencias representa
correctamente las 3 partes de un JWT separadas por puntos?
\begin{itemize}
\item[a)] payload.signature.header
\item[b)] header.signature.payload
\item[c)] header.payload.signature
\item[d)] signature.header.payload
\end{itemize}
\textbf{Respuesta correcta:} (c).

\textbf{Pregunta A2 (V/F):} En una API REST, el método HTTP PUT se
usa para crear un recurso nuevo cuando el identificador aún no existe
en el servidor, mientras que POST actualiza un recurso existente.

\textbf{Respuesta:} \emph{Falso}. PUT reemplaza/actualiza un recurso
identificado por la URL; POST crea un nuevo recurso. Aunque PUT puede
crear si el recurso no existe, su semántica primaria es reemplazar.
\end{cajaejemplo}

\begin{cajaejemplo}{Muestra de la Parte B (respuesta corta)}
\textbf{Pregunta B1:} Explique la diferencia entre sesión y cookie en
el contexto de la gestión de estado en aplicaciones web. Mencione
dónde se almacena cada una.

\textbf{Respuesta esperada:} Una cookie es un fragmento de texto que
se almacena en el navegador del cliente y se envía al servidor en
cada petición. Una sesión es un conjunto de datos asociado a un
usuario que se almacena en el servidor (memoria, Redis, BD); el
cliente solo guarda el ID de sesión en una cookie.
\end{cajaejemplo}

\begin{cajaejemplo}{Muestra de la Parte C (análisis de código)}
\textbf{El siguiente código PHP gestiona el inicio de sesión. Contiene
3 errores. Identifíquelos:}

\begin{lstlisting}[style=php,numbers=left,caption={Ejemplo de PHP},label=lst:17.1]
<?php
session_start(); // arranca o reanuda la sesion
$user = $_GET['user'];
$pwd  = $_GET['pwd'];

$sql = "SELECT * FROM users WHERE name='$user' AND pwd='$pwd'";
$res = mysql_query($sql);

if (mysql_num_rows($res) > 0) { // condicional
    $_SESSION['logueado'] = true;
    setcookie('SESS_REMEMBER', $user, time()+3600);
    echo "Bienvenido $user!"; // imprime al output del servidor
}
\end{lstlisting}

\textbf{Errores esperados:}
\begin{enumerate}
\item Credenciales por GET (deberían ser POST).
\item SQL Injection: concatenación directa en línea 6.
\item \texttt{mysql\_*}: extensión obsoleta y eliminada en PHP 7+;
debe usar PDO o mysqli.
\item Password en texto plano (no se hashea con
\texttt{password\_hash}/\texttt{password\_verify}).
\item Cookie \texttt{SESS\_REMEMBER} sin flags
\texttt{HttpOnly}/\texttt{Secure}/\texttt{SameSite}.
\item XSS en \texttt{echo \char34{}Bienvenido \$user!\char34{}}: falta
\texttt{htmlspecialchars}.
\item No se regenera el ID de sesión tras login (anti-fixation).
\end{enumerate}
(Aunque el enunciado pide 3, hay más; con identificar los 3 más
graves se obtiene puntuación máxima.)
\end{cajaejemplo}

\begin{cajaejemplo}{Muestra de la Parte D (ejercicio integrador)}
\textbf{Implemente en el lenguaje y framework de su elección un
endpoint POST /api/v1/envios que:}

\begin{enumerate}
\item Requiera autenticación con JWT en \texttt{Authorization: Bearer
\dots}.
\item Reciba JSON con \texttt{origen}, \texttt{destino},
\texttt{peso\_kg}.
\item Valide que \texttt{origen} y \texttt{destino} no estén vacíos
y que \texttt{peso\_kg} $>$ 0.
\item Persista en la BD usando ORM o prepared statements.
\item Retorne código HTTP 201 con el envío creado o 422 con
\texttt{errors} en caso de fallo de validación.
\end{enumerate}

Indique al inicio del código el framework y el lenguaje utilizados.
\end{cajaejemplo}

\section{Preparación efectiva para la EP2}

Tener clara la estrategia de preparación durante los siete
días previos al examen y la estrategia durante las dos horas del
mismo es tan importante como el contenido.


\subsection{Plan de estudio de 7 días previos al examen}

La caja siguiente propone un plan día por día para los
siete días anteriores a la EP2.


\begin{cajabuenas}{Plan de 7 días}
\begin{description}[leftmargin=1.5cm]
\item[Día -7:] Re-leer las cajas de definición y buenas prácticas
de las semanas 10--16. Marcar las que no se entienden.
\item[Día -6:] Resolver el ejercicio integrador de las semanas 10
(carrito), 11 (CRUD ORM) y 12 (C4) sin mirar el libro.
\item[Día -5:] Resolver los ejercicios de las semanas 13 (cache),
14 (MVC) y 15 (API REST) sin mirar el libro.
\item[Día -4:] Estudiar el OWASP Top 10:2021 hasta poder explicar
cada uno con un ejemplo.
\item[Día -3:] Practicar análisis de código: tomar 5 fragmentos
con vulnerabilidades y encontrar todos los errores.
\item[Día -2:] Repasar los códigos HTTP y los flags de seguridad de
cookies.
\item[Día -1:] Descansar; revisar solo el resumen ejecutivo de cada
semana. NO empezar a estudiar temas nuevos la noche anterior.
\end{description}
\end{cajabuenas}

\subsection{Estrategia durante el examen}

Durante el examen, la estrategia para distribuir el tiempo
y atacar las preguntas en el orden correcto puede cambiar varios
puntos la calificación final.


\begin{enumerate}
\item Llegar 15 minutos antes con cédula.
\item Leer todo el examen en 5 minutos antes de empezar.
\item Distribución sugerida del tiempo (120 min):
\begin{itemize}
\item Parte A: 15 min.
\item Parte B: 20 min.
\item Parte C: 35 min.
\item Parte D: 45 min.
\item Revisión final: 5 min.
\end{itemize}
\item En la Parte D escribir primero un comentario al inicio:
\texttt{// Framework: Spring Boot 4 / Lenguaje: Java 21}.
\item Si no recuerda exactamente la sintaxis, hacer pseudocódigo
claro: vale más comprensión correcta que sintaxis perfecta.
\item Verificar que cada respuesta de la Parte B tenga máximo 3
líneas; respuestas largas pueden penalizarse.
\end{enumerate}

\section{Reflexión sobre el segundo corte}

Las semanas 10--17 representan la transición del estudiante de
\emph{programador} a \emph{ingeniero de software}. La diferencia es
sutil pero crucial: el programador escribe código que funciona; el
ingeniero documenta decisiones, mide rendimiento, prevé fallos,
considera la seguridad desde el diseño y entiende que un sistema
existe en un contexto humano y organizacional.

El PFC es el primer proyecto de la carrera donde estos elementos
convergen. La defensa oral es un ensayo de las defensas profesionales
que vendrán: revisiones técnicas con clientes, presentaciones a
inversores, defensa de la tesis de grado.

\section{Contraste bibliográfico de los conceptos clave de la EP2}

La EP2 cubre las Unidades III y IV con base en literatura técnica
estándar. La gestión de estado y el acceso a datos están tratados a
profundidad en \cite{fowler2002patterns} y \cite{kleppmann2017designing}.
Los patrones arquitectónicos se apoyan en \cite{bass2021software} y
\cite{richards2020fundamentals}. El modelo C4 está documentado por
\cite{brown2023c4}. REST y SOAP en \cite{fielding2000rest},
\cite{richardson2007restful}, \cite{alonso2010webservice} y
\cite{richardson2007restful}. JWT en RFC 7519 \cite{jones2015jwt} y
\cite{jones2015jwt}. Las pruebas en \cite{martin2008clean} y
\cite{martin2017clean}. La seguridad en \cite{owasp2021top10},
\cite{verizon2025dbir} y \cite{ortega2022securecoding}. El
estudiante debe estar familiarizado con al menos las ideas
principales de cada una de estas fuentes para responder con solvencia
en la EP2 y la defensa oral.

\section{Ejercicio resuelto adicional con resultado visual}

Como cierre, el siguiente ejercicio resuelto adicional
simula una pregunta del tribunal con la respuesta esperada, y se
complementa con un propuesto: simulacro de defensa entre equipos.


\begin{cajaejemplo}{Ejercicio resuelto 17.1 --- Simulación de pregunta-respuesta de defensa}
\textbf{Enunciado.} El tribunal hace una pregunta clásica:
\emph{«¿Por qué no usaron microservicios en su PFC?»} Construir una
respuesta de un minuto que evidencie razonamiento arquitectónico
maduro.

\textbf{Mala respuesta (a evitar):} \emph{«Porque era más fácil con
monolito.»}

\textbf{Buena respuesta:}
\begin{quote}
\small Decidimos arquitectura de monolito modular con base en tres
criterios. \textbf{Primero}, el equipo tiene 4 personas, y según
\cite{newman2019monolith} los microservicios pagan dividendos a
partir de 10--15 desarrolladores; antes de eso, el overhead operativo
supera al beneficio. \textbf{Segundo}, los requerimientos no presentan
necesidades de escalado independiente por módulo: el carrito y el
catálogo crecen al mismo ritmo. \textbf{Tercero}, modularizamos
internamente con paquetes Maven separados
(\texttt{usuarios}, \texttt{productos}, \texttt{pedidos},
\texttt{notificaciones}), siguiendo la idea de
\cite{auer2021monolithic}: monolito hoy, microservicios mañana si
hace falta. Esto nos da despliegue simple en Docker Compose con
posibilidad de migrar si las métricas lo justifican.

La ADR-002 del repositorio documenta esta decisión con criterios y
alternativas consideradas.
\end{quote}

\textbf{Resultado de esta respuesta} (vista del tribunal mirándose
entre sí):

\begin{center}
\fbox{\begin{minipage}{0.85\textwidth}
\vspace{0.15cm}\small
\hspace{0.3cm}\textbf{[Tribunal --- vocal 1]} \emph{«Bien argumentado.»}\\[0.15cm]
\hspace{0.3cm}\textbf{[Tribunal --- vocal 2]} \emph{«¿Y cómo medirían
si en el futuro toca migrar?»}\\[0.15cm]
\hspace{0.3cm}\textbf{[Estudiante]} \emph{«Monitorearíamos métricas
con Prometheus: si el módulo de pedidos consume {$>$}70\% de CPU
mientras otros están al 10\%, sería señal de candidato a extraer
como microservicio.»}\\[0.15cm]
\hspace{0.3cm}\textbf{[Tribunal]} \emph{«Excelente.»}\\
\vspace{0.1cm}
\end{minipage}}
\end{center}

\textbf{Discusión.} La diferencia entre la mala y la buena respuesta
no es la decisión técnica (en ambos casos es monolito), sino el
razonamiento explícito que la acompaña: criterios, autores
referenciados, alternativas, métricas de revisión futura. Esto es lo
que distingue a un \emph{programador} de un \emph{ingeniero de
software} \cite{martin2017clean}.
\end{cajaejemplo}

\begin{cajaejercicio}{Ejercicio propuesto 17.2 --- Simulacro de defensa entre equipos}
\textbf{Enunciado.} Antes de la defensa real, organizar un simulacro
en clase donde cada equipo defiende su PFC ante otros dos equipos que
hacen las preguntas (15--20 minutos por equipo).

\textbf{Procedimiento:}
\begin{enumerate}
\item Cada equipo prepara su presentación de 20 min siguiendo la
estructura recomendada en la sección \emph{«Cómo preparar la defensa
oral»}.
\item Los equipos evaluadores reciben una rúbrica con criterios:
\begin{itemize}
\item Claridad de la exposición técnica (1--5).
\item Justificación de decisiones arquitectónicas (1--5).
\item Calidad de la demo (1--5).
\item Manejo de preguntas (1--5).
\item Distribución equitativa entre integrantes (1--5).
\end{itemize}
\item Tras la defensa, los equipos evaluadores escriben dos
fortalezas y dos mejoras concretas para el equipo evaluado.
\item Cada equipo entrega un acta con las observaciones recibidas y
explica cómo las incorporará antes de la defensa real.
\end{enumerate}

\textbf{Producto final:} dossier del equipo con: rúbrica completada
por cada equipo evaluador, acta de observaciones recibidas, plan de
mejoras priorizado. Subir al repositorio Git como
\texttt{docs/defensa-simulacro.md}.

\textbf{Beneficio pedagógico.} Estudios sobre evaluación entre pares
\cite{washizaki2024swebok} muestran que el simulacro entre equipos
detecta entre 40\% y 60\% de las debilidades de presentación que el
tribunal real detectaría, dando al equipo la oportunidad de
corregirlas. Las preguntas de pares suelen además ser más duras que
las del tribunal académico.
\end{cajaejercicio}

% =====================================================================
%       SEMANA 18: TUTORÍA DE REFUERZO FINAL Y CIERRE DEL CURSO
% =====================================================================
\chapter{Tutoría de refuerzo final y cierre del curso}
\label{cap:semana18}

\epigraph{«Lo que sabemos es una gota; lo que ignoramos, un océano.
Pero la dirección de la gota es lo que importa.»}{--- \textit{Adaptado de
Isaac Newton}}

\section*{Naturaleza de esta semana}
La semana 18 constituye la \textbf{segunda semana de tutoría de
refuerzo} del curso, paralela a la de la semana 8 pero con un alcance
más amplio: cierra el segundo corte y, a la vez, cierra integralmente
la asignatura. Sus contenidos no se planifican rígidamente
\emph{a priori}: se ajustan dinámicamente a las dificultades observadas
en las semanas 10--17 y, particularmente, a los errores frecuentes de
la \emph{Evaluación Parcial Segundo Corte} de la semana 17. Este
capítulo ofrece un guión estructurado de revisión, ejercicios de
consolidación y un cierre profesional del curso.

Los cuatro propósitos de la semana son:

\begin{enumerate}
\item \textbf{Resolver dificultades} en los conceptos donde el grupo
mostró menor desempeño en la EP2.
\item \textbf{Realimentar individualmente} a cada equipo sobre los
PFC entregados, identificando fortalezas y oportunidades de mejora.
\item \textbf{Consolidar} con ejercicios integradores que cubren
transversalmente las cuatro unidades del curso.
\item \textbf{Realizar el cierre integrador} situando el aprendizaje
dentro de la trayectoria profesional del estudiante en el mercado
laboral 2026.
\end{enumerate}

\section{Mapa integrador de la asignatura}

Antes del repaso intensivo, conviene visualizar
globalmente lo que el estudiante ha construido en 17 semanas: cómo
las cuatro unidades se conectan entre sí y convergen en el PFC.


\subsection{Las cuatro unidades como una historia coherente}

A lo largo de 17 semanas hemos transitado un camino que tiene
arquitectura interna: cada unidad se apoya en la anterior y todas
convergen en el Proyecto Fin de Curso. El siguiente diagrama
sintetiza esa estructura:

\begin{figure}[htbp]
\centering
\begin{tikzpicture}[
  every node/.style={font=\scriptsize,align=center},
  unidad/.style={rectangle,draw=uteqverde,thick,fill=uteqverde!10,
                 rounded corners=4pt,minimum width=3.5cm,
                 minimum height=2cm}]

  \node[unidad] (u1) {\textbf{Unidad I}\\\textbf{Cliente}\\
                      HTML5 + CSS3\\+ JavaScript\\Sem 1--4};
  \node[unidad,right=0.4cm of u1] (u2) {\textbf{Unidad II}\\\textbf{Servidor}\\
                                          PERL + PHP\\+ Java + ASP.NET\\Sem 5--9};
  \node[unidad,below=0.4cm of u1] (u3) {\textbf{Unidad III}\\\textbf{Datos y patrones}\\
                                          Estado + ORM\\+ Arquitectura\\Sem 10--13};
  \node[unidad,below=0.4cm of u2] (u4) {\textbf{Unidad IV}\\\textbf{MVC y servicios}\\
                                          MVC + REST/SOAP\\+ Pruebas\\Sem 14--17};

  \node[circle,draw=uteqdorado,thick,fill=uteqdorado!20,
        minimum size=2cm,below right=-0.3cm and 1cm of u2,
        font=\bfseries\scriptsize,align=center]
        (pfc) {PFC\\v1.0.0};

  \draw[->,thick,>=stealth,uteqverde!70!black]
       (u1) -- (u2);
  \draw[->,thick,>=stealth,uteqverde!70!black]
       (u2) -- (u4);
  \draw[->,thick,>=stealth,uteqverde!70!black]
       (u3) -- (u4);
  \draw[->,thick,>=stealth,uteqverde!70!black]
       (u1) -- (u3);

  \draw[->,thick,>=stealth,uteqdorado!70!black,dashed]
       (u4) -- (pfc);
  \draw[->,thick,>=stealth,uteqdorado!70!black,dashed]
       (u3) -- (pfc);

\end{tikzpicture}
\caption{Mapa integrador del curso: las cuatro unidades convergen en
el Proyecto Fin de Curso. Las flechas verdes representan la dependencia
secuencial entre unidades; las flechas doradas, la integración
explícita hacia el PFC.}
\end{figure}

\subsection{Tabla resumen de tecnologías dominadas}

La tabla siguiente lista todas las tecnologías que el
estudiante debe dominar al completar la asignatura, con el nivel
esperado en cada una.


\begin{table}[htbp]
\centering
\caption{Tecnologías que el estudiante debe dominar al completar la asignatura.}
\footnotesize
\begin{tabularx}{\textwidth}{lXl}
\toprule
\textbf{Categoría} & \textbf{Tecnología} & \textbf{Nivel esperado}\\
\midrule
Cliente HTML & HTML5 semántico, ARIA, validación nativa & Avanzado\\
Cliente CSS & CSS3, Flexbox, Grid, mobile-first, Sass & Avanzado\\
Cliente JS & ES2024, DOM, async/await, fetch, módulos & Avanzado\\
Servidor PHP & PHP 8.4, PDO, sesiones, Slim/Laravel & Intermedio\\
Servidor Java & Spring Boot 4, JPA, Spring Security & Avanzado\\
Servidor C\# & ASP.NET Core 10, EF Core & Básico\\
Servidor Perl & CGI, sintaxis básica & Introductorio\\
Datos & PostgreSQL, MySQL, JDBC, JPA, EF Core, ORM & Intermedio\\
Caché & Redis cache-aside con TTL & Intermedio\\
Frontend SPA & Angular 19+, TypeScript, RxJS & Básico\\
Arquitectura & Capas, hexagonal, REST, microservicios & Intermedio\\
Documentación & C4, ADRs, OpenAPI 3.0 & Intermedio\\
Testing & JUnit 5, Mockito, MockMvc, JaCoCo & Intermedio\\
Pruebas de carga & k6 con scripts JavaScript & Básico\\
Frontend audit & Lighthouse, axe DevTools & Intermedio\\
Seguridad & OWASP Top 10:2021, JWT, HTTPS & Intermedio\\
Containerización & Docker, Docker Compose & Intermedio\\
CI/CD & GitHub Actions, pipelines básicos & Básico\\
Versionado & Git, GitHub, Conventional Commits & Avanzado\\
\bottomrule
\end{tabularx}
\end{table}

\section{Errores comunes detectados en las semanas 10--17}

Antes del repaso teórico de los temas, conviene listar los errores
más frecuentes detectados en los entregables de la segunda mitad del
curso y en la EP2. Cada estudiante debe verificar que ya no comete
ninguno de ellos:

\begin{cajaadvertencia}{Diez errores frecuentes que se debe haber superado}
\begin{enumerate}
\item Olvidar las flags \texttt{HttpOnly}, \texttt{Secure} y
\texttt{SameSite} al crear cookies de sesión o de JWT.
\item Concatenar SQL con datos del usuario en lugar de usar parámetros
preparados (PDO, JDBC, EF Core).
\item Confundir HTTP 401 (\emph{Unauthorized} --- no autenticado) con
403 (\emph{Forbidden} --- autenticado pero sin permiso).
\item Usar PUT para crear y POST para actualizar (cuando es al revés).
\item Almacenar JWT en \texttt{localStorage} en lugar de cookie
\texttt{HttpOnly}, exponiéndolo a robo por XSS.
\item Diseñar URIs con verbos como \texttt{/getProductos} o
\texttt{/crearPedido} en lugar de sustantivos como \texttt{/productos}
con métodos HTTP semánticos.
\item Olvidar el claim \texttt{exp} en el JWT, generando tokens que
nunca expiran.
\item Saltar a microservicios cuando un monolito modular bien
estructurado bastaría (regla \emph{«monolith first»} de Fowler).
\item No invalidar la caché al modificar la BD, generando
inconsistencias visibles para el usuario.
\item Documentar la API solo con un README en lugar de exponer
OpenAPI 3.0 con Swagger UI interactivo.
\end{enumerate}
\end{cajaadvertencia}

\section{Repaso intensivo de los temas más difíciles}

Basado en el análisis estadístico de errores en la EP2, los siguientes
cuatro temas concentran la mayor parte de las pérdidas de puntaje y
requieren refuerzo específico.

\subsection{Tema 1: Diferencia REST vs SOAP}

La diferencia entre REST y SOAP es una de las preguntas más
frecuentes del tribunal y de la EP2. El resumen siguiente captura lo
esencial.


\begin{cajadefinicion}{Resumen ejecutivo --- REST vs SOAP}
\textbf{REST} es un \emph{estilo arquitectónico} sobre HTTP: usa
verbos HTTP semánticamente, recursos como sustantivos en URLs, JSON
como formato dominante. Es \emph{stateless} por restricción.

\textbf{SOAP} es un \emph{protocolo} sobre HTTP/SMTP/otros: define
mensajes XML con \texttt{envelope} y \texttt{body}, contrato formal
con WSDL, soporta WS-* para seguridad y transacciones.

\textbf{Cuándo elegir cada uno}: REST para todo lo nuevo, especialmente
APIs públicas y SPAs. SOAP solo si hay obligación de integrar con
sistemas legados (banca, salud, B2B grande).
\end{cajadefinicion}

\subsection{Tema 2: JWT --- estructura y uso correcto}

JWT aparece en toda API REST moderna del curso. Conocer su
estructura exacta y las reglas de uso correcto es prerrequisito para
defender el PFC con solvencia.


\begin{cajadefinicion}{Resumen ejecutivo --- JWT}
Un JWT tiene \textbf{tres partes separadas por puntos}:
\textbf{header.payload.signature}.

\begin{itemize}
\item Header: algoritmo de firma (HS256, RS256).
\item Payload: claims (sub, iss, exp, iat, jti, custom).
\item Signature: HMAC o RSA del header+payload con secreto.
\end{itemize}

\textbf{Reglas de uso correcto:}
\begin{itemize}
\item Almacenar en cookie \texttt{HttpOnly} \texttt{Secure}
\texttt{SameSite=Strict} (NO en localStorage).
\item Tiempo de expiración corto (1 hora típico).
\item Refresh token separado con expiración mayor (7 días).
\item Blacklist en Redis por JTI para revocación.
\item Validar firma en cada petición.
\end{itemize}
\end{cajadefinicion}

\subsection{Tema 3: ORM y problema N+1}

El problema N+1 sigue siendo la patología de rendimiento
más frecuente en aplicaciones que usan ORM. La caja siguiente lo
sintetiza junto con su solución.


\begin{cajadefinicion}{Resumen ejecutivo --- ORM y N+1}
El ORM mapea objetos a tablas. La trampa principal es el \textbf{N+1}:
al iterar una colección con relaciones, el ORM ejecuta una consulta
adicional por cada elemento.

\textbf{Solución}: \emph{eager loading} explícito.
\begin{itemize}
\item Spring Data: \texttt{@Query(\char34{}SELECT p FROM Producto p JOIN
FETCH p.categoria\char34{})}.
\item EF Core: \texttt{.Include(p => p.Categoria)}.
\item Eloquent: \texttt{Producto::with('categoria')->get()}.
\item Doctrine: \texttt{SELECT p, c FROM Producto p JOIN p.categoria c}.
\end{itemize}
\end{cajadefinicion}

\subsection{Tema 4: Cache-aside con Redis}

Cache-aside con Redis es el patrón de caché dominante en
2026. La caja siguiente sintetiza su mecánica y los TTL típicos.


\begin{cajadefinicion}{Resumen ejecutivo --- cache-aside}
El patrón cache-aside funciona así:

\begin{enumerate}
\item App consulta Redis con la \emph{key}.
\item Si \textbf{HIT}: devuelve valor cacheado.
\item Si \textbf{MISS}: consulta BD, guarda en Redis con TTL,
devuelve.
\end{enumerate}

\textbf{Al actualizar la BD}: \emph{invalidar} el cache (borrar la
key), no actualizar la cache directamente. Esto evita inconsistencia
si la operación de BD falla a mitad.

\textbf{TTL típicos:} 60s para listados frecuentes; 5--15 min para
catálogos relativamente estables; 1 hora para datos casi inmutables.
\end{cajadefinicion}

\section{Banco de ejercicios de consolidación}

Los siguientes tres ejercicios integran transversalmente las
competencias adquiridas en las semanas 10--17. Se sugiere al
estudiante elegir al menos uno e implementarlo durante esta semana
de cierre como ejercicio voluntario de consolidación.

\begin{cajaejercicio}{Ejercicio 18.1 --- API REST integral de logística}
\textbf{Objetivo:} diseñar e implementar una API REST completa de
gestión de envíos para una empresa de logística, integrando todos los
elementos clave del segundo corte.

\textbf{IDE:} IntelliJ IDEA Ultimate; \textbf{stack}: Spring Boot 4
+ PostgreSQL 16 + Redis 7.

\textbf{Requisitos:}
\begin{itemize}
\item Autenticación JWT con \texttt{accessToken} (1 h) y
\texttt{refreshToken} (7 días) en cookies HttpOnly.
\item Endpoints: \texttt{POST /api/v1/envios},
\texttt{GET /api/v1/envios/\{id\}},
\texttt{GET /api/v1/envios?estado=...\&page=...},
\texttt{PATCH /api/v1/envios/\{id\}/estado}.
\item Validación con \texttt{@Valid} y Bean Validation.
\item Cache Redis con TTL de 60 s en el listado de envíos.
\item Documentación OpenAPI 3.0 generada automáticamente con
springdoc.
\item Manejo de errores con ProblemDetails (RFC 7807).
\item Cobertura JaCoCo $\geq 70\%$.
\item Prueba con Postman: 12 escenarios incluyendo casos de error
(400, 401, 403, 404, 422, 429).
\end{itemize}
\end{cajaejercicio}

\begin{cajaejercicio}{Ejercicio 18.2 --- Migración SOAP a REST}
\textbf{Objetivo:} ejercitar la integración con sistemas legados
y la justificación arquitectónica documentada.

\textbf{Procedimiento:}
\begin{enumerate}
\item Tomar el WSDL de un servicio SOAP heredado (proporcionado por
el docente; alternativamente, un servicio público como el de tipos de
cambio del Banco Central del Ecuador).
\item Construir una capa de adaptación REST en Spring Boot que
exponga los métodos SOAP como recursos REST con JSON.
\item Implementar caché Redis para reducir la carga sobre el servicio
SOAP origen.
\item Producir una ADR (\emph{Architectural Decision Record}) que
justifique la decisión: por qué se mantiene el servicio SOAP
en lugar de reescribirlo, qué se gana al ofrecer la fachada REST,
qué riesgos asume el equipo.
\item Pruebas Postman comparando latencia con cache fría vs caliente.
\end{enumerate}
\end{cajaejercicio}

\begin{cajaejercicio}{Ejercicio 18.3 --- Auditoría de seguridad del PFC}
\textbf{Objetivo:} aplicar OWASP ZAP en modo \emph{baseline} sobre el
PFC entregado y producir un informe profesional de vulnerabilidades.

\textbf{Procedimiento:}
\begin{enumerate}
\item Ejecutar el PFC localmente con \texttt{docker compose up}.
\item Descargar OWASP ZAP desde \url{https://www.zaproxy.org}.
\item Configurar ZAP en modo \emph{Automated Scan} apuntando a la URL
del frontend Angular y la URL de la API REST.
\item Ejecutar el escaneo (10--30 min según el tamaño del sistema).
\item Producir un informe que incluya:
\begin{itemize}
\item Tabla de vulnerabilidades con severidad (Alta, Media, Baja,
Informativa).
\item Por cada vulnerabilidad: descripción, impacto, recomendación
concreta de corrección y referencia OWASP.
\item Plan de remediación priorizado.
\end{itemize}
\item Subsanar al menos las vulnerabilidades de severidad alta y
re-ejecutar el escaneo para evidenciar el cierre.
\end{enumerate}
\end{cajaejercicio}

\section{Realimentación individual del PFC}

Tras la presentación grupal en la semana 17, cada equipo
recibe realimentación individualizada en la semana 18 sobre su
PFC: qué hicieron muy bien, qué mejorarían.


\subsection{Criterios de excelencia observados}

Los PFC que alcanzaron las calificaciones más altas comparten estas
características:

\begin{itemize}
\item \textbf{Decisiones documentadas}: cada decisión técnica
significativa tiene una ADR justificando la elección.
\item \textbf{Pruebas reales}: cobertura $\geq$ 80\% con tests
significativos, no triviales.
\item \textbf{Demo robusta}: \texttt{docker compose up} funciona en
una computadora limpia, sin trucos.
\item \textbf{Métricas medidas}: scores de Lighthouse y resultados de
k6 reales, no inventados.
\item \textbf{Equipo cohesionado}: cualquier integrante puede
responder cualquier pregunta del tribunal.
\end{itemize}

\subsection{Mejoras frecuentemente sugeridas}

Las observaciones más frecuentes que reciben los equipos
en esta realimentación son recomendaciones de evolución futura,
no defectos.


\begin{itemize}
\item Agregar pruebas de mutación con PIT (Pitest) para evaluar la
calidad real de las pruebas, no solo su cobertura.
\item Configurar autoscaling con Kubernetes en lugar de Docker
Compose para entornos productivos.
\item Implementar OAuth 2.0 / OpenID Connect en lugar de JWT casero
para integración con identidad institucional.
\item Agregar observabilidad: Prometheus + Grafana + Tempo para
métricas, dashboards y traces distribuidos.
\item Implementar feature flags para despliegues progresivos sin
re-build (LaunchDarkly, Unleash).
\item Migrar el frontend a SSR con Angular Universal para mejorar
SEO y la métrica LCP.
\end{itemize}

\section{El curso en perspectiva profesional}

El curso termina, pero la formación profesional sigue. La
sección siguiente sitúa lo aprendido en la perspectiva de la
carrera del estudiante en el mercado laboral 2026.


\subsection{Cómo se ve el desarrollo web profesional en 2026}

El campo del desarrollo web ha alcanzado una madurez que hace 10
años parecía lejana. En 2026:

\begin{itemize}
\item El desarrollador full-stack típico domina al menos un framework
backend (Spring Boot, ASP.NET Core, Node.js, Laravel) y uno frontend
(Angular, React, Vue).
\item Containers (Docker) y Kubernetes son habilidades esperadas.
\item TypeScript ha desplazado a JavaScript en proyectos serios.
\item La IA generativa (GitHub Copilot, Claude, ChatGPT) duplicó la
productividad declarada del desarrollador medio.
\item La seguridad ya no es opcional: se exigen auditorías OWASP en
contratos.
\item Los equipos remotos son la norma, con asincronía como cultura.
\end{itemize}

\subsection{Próximos pasos en la formación}

El siguiente paso natural en la malla curricular es la asignatura de
\emph{Despliegue y Operaciones} de sexto nivel, donde se profundiza
en CI/CD, contenedores avanzados, observabilidad y prácticas DevOps.
Adicionalmente, recomendamos al estudiante:

\begin{itemize}
\item \textbf{Profundizar en un framework}: elegir uno (Spring Boot,
.NET, Laravel) y dominarlo a nivel experto.
\item \textbf{Aprender un orquestador}: Kubernetes es la habilidad
más demandada en infraestructura.
\item \textbf{Bases de datos avanzadas}: índices, optimización,
sharding, replicación.
\item \textbf{Diseño de sistemas distribuidos}: leer
\emph{Designing Data-Intensive Applications} de Martin Kleppmann.
\item \textbf{Patrones de microservicios}: Sam Newman, Chris
Richardson.
\item \textbf{Seguridad ofensiva}: TryHackMe, HackTheBox,
certificaciones (eJPT, OSCP).
\item \textbf{Inglés técnico}: la mayor parte del conocimiento
publicado y los empleos remotos están en inglés.
\end{itemize}

\subsection{Certificaciones que aportan valor profesional en 2026}

Para complementar la formación con credenciales reconocidas
por el mercado, hay certificaciones de alto valor profesional que el
estudiante puede considerar tras completar el curso.


\begin{table}[htbp]
\centering
\caption{Certificaciones de alto valor profesional para el desarrollador web 2026.}
\footnotesize
\begin{tabularx}{\textwidth}{lXl}
\toprule
\textbf{Categoría} & \textbf{Certificación} & \textbf{Valor mercado}\\
\midrule
Cloud & AWS Certified Developer Associate & Alto\\
Cloud & Microsoft Azure Developer (AZ-204) & Alto\\
Cloud & Google Professional Cloud Developer & Alto\\
Containers & Certified Kubernetes Application Developer (CKAD) & Muy alto\\
Java & Oracle Certified Professional Java SE 21 & Medio\\
.NET & MS Certified: .NET Developer & Medio\\
Spring & VMware Spring Professional & Alto\\
Seguridad & CompTIA Security+ & Alto\\
Seguridad & Certified Ethical Hacker (CEH) & Medio\\
Datos & MongoDB Certified Developer & Medio\\
\bottomrule
\end{tabularx}
\end{table}

\section{Conclusión del curso}

El curso \emph{Aplicaciones Web} ha cubierto el ciclo completo del
desarrollo web profesional: desde el HTML semántico que un navegador
renderiza hasta la arquitectura distribuida con caché Redis y
balanceadores Nginx que sirve a miles de usuarios. El estudiante que
ha completado las 18 semanas posee:

\begin{itemize}
\item Capacidad para construir interfaces web semánticas, accesibles
y responsivas conforme al W3C y WCAG 2.1.
\item Dominio operativo de cuatro lenguajes y plataformas servidor:
PERL/CGI, PHP, Java/Spring Boot y ASP.NET Core.
\item Diseño e implementación de APIs REST documentadas con OpenAPI,
autenticadas con JWT y respaldadas por bases de datos relacionales
con acceso vía ORM.
\item Razonamiento arquitectónico: capas, MVC, microservicios, EDA,
P2P y modelo C4.
\item Conciencia de seguridad: OWASP Top 10, controles defensivos,
cifrado, gestión de identidad.
\item Experiencia de proyecto fin de curso \emph{end-to-end}: desde
la planificación (Entrega 1A en semana 5) hasta el despliegue final
(semana 16) y la defensa oral (semana 17).
\end{itemize}

Las decisiones técnicas que el estudiante ha documentado, los patrones
que ha aplicado y las pruebas que ha escrito en su PFC no son
ejercicios académicos sin más: son \emph{habilidades transferibles
directamente al mercado laboral}. Las empresas locales ---desde
\emph{startups} de Quevedo hasta empresas medianas de Guayaquil y
Quito--- buscan precisamente este perfil. Las plataformas de empleo
remoto (Toptal, Turing, Crossover, Andela LATAM) abren además puertas
internacionales para quienes complementen este conocimiento técnico
con inglés y disciplina profesional.

El docente felicita a cada estudiante por completar el camino y
recuerda que el aprendizaje del desarrollo de software no termina
nunca: empieza otra vez cada lunes con un nuevo \emph{commit}.

\hfill \emph{Buen viaje. Y buen código.}\\
\hfill \emph{Mayo de 2026, Quevedo.}

\section{Cierre administrativo de la asignatura}

Los siguientes puntos cierran administrativamente la
asignatura conforme a la Norma Académica UTEQ.


\begin{itemize}
\item \textbf{Calificación final} = ponderación según Norma Académica
UTEQ de los componentes GA, TA y EV.
\item \textbf{Aprobación}: $\geq$ 7.0 en escala 0--10.
\item \textbf{Recuperación}: examen integral cubriendo las 4 unidades
para estudiantes con calificación final entre 5.0 y 6.9.
\item \textbf{Repetición}: estudiantes con calificación $<$ 5.0 deben
recursar la asignatura el siguiente período.
\item \textbf{Acta de notas}: se publica oficialmente 7 días hábiles
después del cierre de la semana 18.
\end{itemize}

% Bibliografia desde archivo externo .bib
% =====================================================================
% APENDICE A - Practica integradora de Unidad II (insertado)
% =====================================================================
% =====================================================================
%  APÉNDICE — PRÁCTICA INTEGRADORA DE UNIDAD II
%  CRUD con patrón DAO «sin framework» en JSP/JSTL/EL, Django y Laravel
%  (Reutiliza el preámbulo del libro. Requiere \lstdefinestyle{python}.)
% =====================================================================

